\begin{chapterpage}{Pengantar data}
  \chaptertitle{Pengantar data}
  \label{introductionToData}
  \label{ch_intro_to_data}
  \chaptersection{basicExampleOfStentsAndStrokes}
  \chaptersection{dataBasics}
  \chaptersection{overviewOfDataCollectionPrinciples}
  % \chaptersection{section_obs_data_sampling}
  \chaptersection{experimentsSection}
\end{chapterpage}
\renewcommand{\chapterfolder}{ch_intro_to_data}

%\begin{tipBox}{\tipBoxTitle[Tujuan Bab:]{Memikirkan data}
%Pahami dasar-dasar organisasi data, jenis data, ringkasan numerik data, ringkasan grafis data, dan teknik dasar pengumpulan data. Kita mengawali dan mengakhiri bab ini dengan studi kasus.}
%\end{tipBox}

\chapterintro{Para ilmuwan berupaya menjawab pertanyaan
  dengan menggunakan metode yang ketat dan pengamatan yang cermat.
  Pengamatan ini -- yang dihimpun, antara lain, dari catatan lapangan,
  survei, dan eksperimen -- menjadi landasan penyelidikan statistika
  dan disebut \term{data}.
  Statistika adalah ilmu tentang cara terbaik untuk mengumpulkan dan menganalisis
  data serta menarik kesimpulan darinya, %Statistika dapat ditempatkan dalam konteks proses penyelidikan secara umum:
%\begin{enumerate}
%\setlength{\itemsep}{0mm}
%\item Tentukan pertanyaan atau masalah.
%\item Kumpulkan data yang relevan dengan topik tersebut.
%\item Analisis data tersebut.
%\item Tarik kesimpulan.
%%\item Ambil keputusan berdasarkan kesimpulan tersebut.
%\end{enumerate}
%Sebagai bidang ilmu, statistika berfokus untuk menjadikan tahap 2--4 objektif, ketat, dan efisien. Dengan~kata lain, statistika memiliki tiga komponen utama: Bagaimana cara terbaik mengumpulkan data? Bagaimana data seharusnya dianalisis? Dan apa yang dapat kita simpulkan dari analisis tersebut?
  dan pada bab pertama ini,
  kita berfokus pada sifat-sifat data
  serta cara pengumpulannya.}

%Topik yang diselidiki para ilmuwan sama beragamnya dengan pertanyaan yang mereka ajukan. Namun, banyak penyelidikan ini dapat ditangani dengan sejumlah kecil teknik pengumpulan data, alat analisis, dan konsep mendasar dalam inferensi statistik. Bab ini memberikan gambaran sekilas tentang tema-tema tersebut dan tema lain yang akan kita jumpai di sepanjang buku. Kita memperkenalkan prinsip dasar setiap cabang dan mempelajari beberapa alat sambil berjalan. Kita akan menjumpai penerapan dari bidang lain, beberapa di antaranya biasanya tidak dikaitkan dengan sains tetapi tetap dapat memperoleh manfaat dari kajian statistik.



\section{Studi kasus: penggunaan stent untuk mencegah strok}
\label{basicExampleOfStentsAndStrokes}

\index{data!strok|(}

Bagian~\ref{basicExampleOfStentsAndStrokes} memperkenalkan tantangan klasik dalam statistika: mengevaluasi efektivitas suatu penanganan medis. Istilah-istilah dalam bagian ini, dan bahkan sebagian besar isi bab ini, akan dibahas kembali nanti. Untuk saat ini, tujuannya hanyalah memperoleh gambaran tentang peran statistika dalam praktik.

Dalam bagian ini kita akan meninjau sebuah eksperimen yang mengkaji efektivitas stent dalam menangani pasien yang berisiko mengalami strok.
Stent adalah alat yang ditempatkan di dalam pembuluh darah untuk membantu pemulihan pasien setelah kejadian kardiak serta mengurangi risiko serangan jantung berikutnya atau kematian. Banyak dokter berharap stent akan memberikan manfaat serupa bagi pasien yang berisiko mengalami strok. Kita mulai dengan menuliskan pertanyaan utama yang hendak dijawab oleh para peneliti:
\begin{quote}
Apakah penggunaan stent mengurangi risiko strok?
\end{quote}

Para peneliti tersebut melakukan eksperimen terhadap 451 pasien berisiko. Setiap pasien sukarelawan dialokasikan secara acak ke salah satu dari dua kelompok:
\begin{itemize}
\item[]\termsub{Kelompok perlakuan}{kelompok perlakuan}. Pasien dalam kelompok perlakuan menerima stent dan penatalaksanaan medis. Penatalaksanaan medis tersebut mencakup pemberian obat, pengelolaan faktor risiko, dan bantuan untuk mengubah gaya hidup.
\item[]\termsub{Kelompok kontrol}{kelompok kontrol}. Pasien dalam kelompok kontrol menerima penatalaksanaan medis yang sama seperti kelompok perlakuan, tetapi tidak menerima stent.
\end{itemize}
Para peneliti mengalokasikan secara acak 224 pasien ke kelompok perlakuan dan 227 pasien ke kelompok kontrol. Dalam penelitian ini, kelompok kontrol menjadi titik acuan untuk mengukur dampak medis stent pada kelompok perlakuan.

Para peneliti mempelajari efek stent pada dua waktu: 30~hari setelah pendaftaran dan 365~hari setelah pendaftaran. Hasil untuk 5 pasien dirangkum dalam Gambar~\ref{stentStudyResultsDF}. Luaran pasien dicatat sebagai ``strok'' atau ``tidak ada kejadian'', yang menunjukkan apakah pasien mengalami strok hingga akhir suatu periode waktu.

\begin{figure}[h]
\centering
\begin{tabular}{l ccc}
\hline
Pasien	&	kelompok	&	0--30 hari 	&	0--365 hari \\
\hline
1		&	perlakuan &	tidak ada kejadian &	tidak ada kejadian \\
2		&	perlakuan &	strok & strok \\
3		&	perlakuan &	tidak ada kejadian & tidak ada kejadian \\
$\vdots$	&	$\vdots$	  &	$\vdots$ \\
450	&	kontrol &	tidak ada kejadian &	tidak ada kejadian \\
451	&	kontrol &	tidak ada kejadian &	tidak ada kejadian \\
\hline
\end{tabular}
\caption{Hasil untuk lima pasien dari penelitian stent.}
\label{stentStudyResultsDF}
% trmt <- c(rep('trmt', 224), rep('control', 227)); outcome30 <- c(rep(c('event', 'no_event'), c(33, 191)), rep(c('event', 'no_event'), c(13, 214))); outcome365 <- c(rep(c('event', 'no_event'), c(33, 191)), rep(c('event', 'no_event'), c(13, 214)))
\end{figure}

Menelaah data setiap pasien satu per satu merupakan cara yang panjang dan rumit untuk menjawab pertanyaan penelitian semula. Sebaliknya, analisis data statistik memungkinkan kita menelaah seluruh data sekaligus. Gambar~\ref{stentStudyResults} merangkum data mentah dengan cara yang lebih berguna. Dari tabel ini, kita dapat dengan cepat melihat apa yang terjadi selama penelitian. Misalnya, untuk menentukan jumlah pasien dalam kelompok perlakuan yang mengalami strok dalam 30 hari, kita melihat sisi kiri tabel pada perpotongan baris perlakuan dan kolom strok: 33.

\begin{figure}[h]
\centering
\begin{tabular}{l cc c cc}
& \multicolumn{2}{c}{0--30 hari} &\hspace{5mm}\ & \multicolumn{2}{c}{0--365 hari} \\
  \cline{2-3} \cline{5-6}
	& 	strok 	& tidak ada kejadian && 	strok 	& tidak ada kejadian \\
  \hline
perlakuan 	& 33		& 191	&&	45 	& 179 \\
kontrol 		& 13		& 214	&& 	28	& 199 \\
  \hline
Jumlah				& 46		& 405	&&	73	& 378 \\
  \hline
\end{tabular}
\caption{Statistik deskriptif untuk penelitian stent.}
\label{stentStudyResults}
\end{figure}

\begin{exercisewrap}
\begin{nexercise}
Dari 224 pasien dalam kelompok perlakuan, 45 orang mengalami strok hingga akhir tahun pertama. Dengan menggunakan kedua bilangan ini, hitung proporsi pasien dalam kelompok perlakuan yang mengalami strok hingga akhir tahun pertama. (Perhatikan: jawaban untuk semua soal Latihan Terbimbing diberikan dalam catatan kaki.)\footnotemark
\end{nexercise}
\end{exercisewrap}\footnotetext{Proporsi dari 224 pasien yang mengalami strok dalam 365 hari: $45/224 = 0.20$.}

Kita dapat menghitung statistik ringkasan dari tabel tersebut.
Sebuah \term{statistik ringkasan}%
\index{statistik|seealso{statistik ringkasan}}
adalah satu bilangan yang merangkum
sejumlah besar data.
Sebagai contoh, hasil utama penelitian setelah 1~tahun
dapat dijelaskan dengan dua statistik ringkasan:
proporsi orang yang mengalami strok dalam kelompok perlakuan
dan kelompok kontrol.
\begin{itemize}
\setlength{\itemsep}{0mm}
\item[] Proporsi yang mengalami strok dalam kelompok perlakuan (stent): $45/224 = 0.20 = 20\%$.
\item[] Proporsi yang mengalami strok dalam kelompok kontrol: $28/227 = 0.12 = 12\%$.
\end{itemize}
Dua statistik ringkasan ini berguna untuk mencari perbedaan antarkelompok, dan kita menemukan sesuatu yang mengejutkan: terdapat tambahan 8\% pasien yang mengalami strok dalam kelompok perlakuan! Hal ini penting karena dua alasan. Pertama, hasil tersebut bertentangan dengan harapan dokter bahwa stent akan \emph{menurunkan} angka kejadian strok. Kedua, hasil tersebut memunculkan pertanyaan statistik: apakah data menunjukkan perbedaan yang ``nyata'' antara kedua kelompok?

Pertanyaan kedua ini cukup rumit. Misalkan Anda melempar koin 100 kali. Walaupun peluang koin menunjukkan sisi kepala pada setiap lemparan adalah 50\%, kemungkinan besar kita tidak akan mengamati tepat 50 sisi kepala. Fluktuasi semacam ini merupakan bagian dari hampir semua proses pembangkitan data. Perbedaan 8\% dalam penelitian stent mungkin saja disebabkan oleh variasi alami ini. Namun, semakin besar perbedaan yang kita amati (untuk ukuran sampel tertentu), semakin sulit dipercaya bahwa perbedaan tersebut terjadi karena kebetulan. Jadi, pertanyaan yang sebenarnya kita ajukan adalah: apakah perbedaannya sedemikian besar sehingga kita harus menolak anggapan bahwa perbedaan tersebut terjadi karena kebetulan?

Meskipun kita belum memiliki perangkat statistik untuk menjawab pertanyaan ini sepenuhnya secara mandiri, kita dapat memahami kesimpulan analisis yang telah dipublikasikan: terdapat bukti kuat bahwa penggunaan stent membahayakan pasien dalam penelitian ini.

\textbf{Berhati-hatilah:}
Jangan menggeneralisasi hasil penelitian ini kepada semua pasien
dan semua jenis stent.
Penelitian ini mengamati pasien dengan karakteristik yang sangat khusus,
yang secara sukarela ikut serta dan mungkin tidak
mewakili semua pasien strok.
Selain itu, terdapat banyak jenis stent, sedangkan penelitian ini hanya
mempertimbangkan stent Wingspan yang dapat mengembang sendiri (Boston Scientific).
Namun, penelitian ini tetap memberi kita pelajaran penting:
kita harus selalu siap menghadapi temuan yang mengejutkan.

\index{data!strok|)}


{\exercisesheader{}

% 1

\eoce{\qt{Migrain dan akupunktur,
    Bagian I\label{migraine_and_acupuncture_intro}}
Migrain merupakan jenis sakit kepala yang sangat nyeri,
dan sebagian pasien berupaya mengatasinya dengan akupunktur.
Untuk menentukan apakah akupunktur meredakan nyeri
migrain, para peneliti melakukan uji terkontrol acak
yang melibatkan 89 perempuan yang didiagnosis menderita migrain.
Mereka dialokasikan secara acak ke salah satu dari dua kelompok:
perlakuan atau kontrol.
Sebanyak 43 pasien dalam kelompok perlakuan menerima akupunktur
yang dirancang secara khusus untuk menangani migrain.
Sebanyak 46 pasien dalam kelompok kontrol menerima akupunktur plasebo
(penusukan jarum pada lokasi yang bukan titik akupunktur).
24 jam setelah menerima akupunktur, pasien ditanya
apakah mereka bebas nyeri.
Hasilnya dirangkum dalam tabel kontingensi
di bawah ini.\footfullcite{Allais:2011}

\noindent\begin{minipage}[l]{0.4\textwidth}
\begin{tabular}{ll  cc c} 
			                         		&           & \multicolumn{2}{c}{\textit{Bebas nyeri}} \\
\cline{3-4}
			                        	 	&			& Ya 	& Tidak 	                  & Jumlah \\
\cline{2-5}
							& Perlakuan 	& 10	 	& 33		                  & 43 \\
\raisebox{1.5ex}[0pt]{\emph{Kelompok}} & Kontrol	 	& 2	 	& 44 	 	                  & 46 \\
\cline{2-5}
							& Jumlah		& 12		& 77		                  & 89
\end{tabular}
\end{minipage}

\begin{quote}
\small\textit{Deskripsi tekstual: Pada skema telinga dalam artikel asli,
titik M berada di dekat bagian bawah depan daun telinga, sedangkan titik S
berada di bagian tengah atas. Artikel tersebut membedakan area akupunktur
yang sesuai (M) dan yang tidak sesuai (S) untuk menangani serangan migrain.
Gambar pihak ketiga tidak disertakan dalam edisi turunan ini.}
\end{quote}
\begin{parts}
\item Berapa persen pasien dalam kelompok perlakuan yang bebas nyeri 24 jam
setelah menerima akupunktur?
\item Berapa persen pasien dalam kelompok kontrol yang bebas nyeri?
\item Kelompok mana yang memiliki persentase pasien bebas nyeri lebih tinggi 24 jam
setelah menerima akupunktur?
\item Temuan Anda sejauh ini mungkin menunjukkan bahwa akupunktur merupakan terapi
yang efektif bagi semua orang yang menderita migrain. Namun, ini bukan satu-satunya
kesimpulan yang dapat ditarik berdasarkan temuan Anda sejauh ini. Apa satu
penjelasan lain yang mungkin atas perbedaan teramati antara persentase
pasien yang bebas nyeri 24 jam setelah menerima akupunktur pada kedua kelompok?
\end{parts}
}{}

% 2

\eoce{\qt{Sinusitis dan antibiotik,
    Bagian I\label{sinusitis_and_antibiotics_intro}} 
Para peneliti yang membandingkan efek pengobatan antibiotik untuk sinusitis akut
dengan penanganan simtomatik mengalokasikan secara acak 166 orang dewasa yang didiagnosis
menderita sinusitis akut ke salah satu dari dua kelompok: perlakuan atau kontrol. Peserta
penelitian menerima amoksisilin (antibiotik) selama 10 hari
atau plasebo yang rupa dan rasanya serupa. Plasebo tersebut terdiri atas
penanganan simtomatik seperti asetaminofen, dekongestan hidung, dan sebagainya.
Pada akhir periode 10 hari, pasien ditanya apakah
gejala mereka membaik.
Distribusi jawaban dirangkum di bawah ini. 
\footfullcite{Garbutt:2012}
\begin{center}
\begin{tabular}{ll  cc c} 
                                    			&			& \multicolumn{2}{c}{\textit{Perbaikan gejala}} \\
                                    			&			& \multicolumn{2}{c}{\textit{menurut laporan pasien}} \\
\cline{3-4}
			                        		&			& Ya 	& Tidak 	& Jumlah \\
\cline{2-5}
							& Perlakuan 	& 66		& 19		& 85 \\
\raisebox{1.5ex}[0pt]{\emph{Kelompok}}	& Kontrol		& 65		& 16 		& 81 \\
\cline{2-5}
							& Jumlah		& 131	& 35		& 166
\end{tabular}
\end{center}
\begin{parts}
\item Berapa persen pasien dalam kelompok perlakuan yang mengalami perbaikan
gejala?
\item Berapa persen pasien dalam kelompok kontrol yang mengalami perbaikan
gejala?
\item Kelompok mana yang memiliki persentase pasien dengan perbaikan
gejala lebih tinggi?
\item
    Temuan Anda sejauh ini mungkin menunjukkan adanya perbedaan nyata
    antara efektivitas antibiotik dan plasebo
    dalam memperbaiki gejala sinusitis.
    Namun, ini bukan satu-satunya kesimpulan yang
    dapat ditarik berdasarkan temuan Anda sejauh ini.
    Apa satu penjelasan lain yang mungkin atas perbedaan teramati
    antara persentase pasien dalam kelompok perlakuan
    antibiotik dan plasebo yang mengalami
    perbaikan gejala sinusitis?
\end{parts}
}{}
}



\section{Dasar-dasar data}
\label{dataBasics}

Pengorganisasian dan pendeskripsian data secara efektif merupakan langkah
awal dalam sebagian besar analisis.
Bagian ini memperkenalkan \emph{matriks data} untuk mengorganisasi data,
serta sejumlah istilah bagi berbagai bentuk data yang akan digunakan
di sepanjang buku ini.

\subsection{Observasi, variabel, dan matriks data}

\index{data!loan50|(}

Gambar~\ref{loan50DF} menampilkan baris 1, 2, 3, dan 50 dari himpunan data
50 pinjaman yang diambil sebagai sampel acak dari pinjaman yang ditawarkan
melalui Lending Club, sebuah perusahaan pinjaman antarpengguna.
Observasi-observasi ini selanjutnya disebut sebagai himpunan data
\data{loan50}.

Setiap baris pada tabel mewakili satu pinjaman.
Nama formal bagi sebuah baris adalah \term{kasus}
atau \term{unit observasi}\index{unit observasi}.
Kolom-kolomnya mewakili karakteristik setiap pinjaman,
yang disebut \termsub{variabel}{variabel}.
Sebagai contoh, baris pertama mewakili pinjaman sebesar \$22,000 dengan tingkat bunga 10.90\%; peminjamnya berdomisili di New Jersey (NJ) dan berpendapatan \$59,000.

\begin{exercisewrap}
\begin{nexercise}
Apa peringkat pinjaman pertama pada Gambar~\ref{loan50DF}?
Dan bagaimana status kepemilikan rumah peminjam
untuk pinjaman pertama tersebut?
Untuk pertanyaan Latihan Terbimbing ini, Anda dapat memeriksa jawaban
di catatan kaki.\footnotemark{}
\end{nexercise}
\end{exercisewrap}
\footnotetext{Peringkat pinjaman tersebut adalah B,
  dan peminjam menyewa tempat tinggalnya.}

Dalam praktik, mengajukan pertanyaan klarifikasi sangatlah penting
untuk memastikan bahwa aspek-aspek penting data telah dipahami.
Misalnya, kita harus selalu memastikan bahwa kita mengetahui makna
setiap variabel dan satuan pengukurannya.
Deskripsi variabel-variabel \data{loan50} diberikan
pada Gambar~\ref{loan50Variables}.

\begin{figure}[h]
\centering
{\small
\begin{tabular}{ccc ccc cc} %c}
  \hline
   & \var{loan\us{}amount}
   & \var{interest\us{}rate}
   & \var{term} & \var{grade} & \var{state}
   & \var{total\us{}income}
   & \var{homeownership} \\
  \hline
  1 & 22000 & 10.90 & 60.00 & B & NJ & 59000.00 & rent \\ 
  2 & 6000 & 9.92 & 36.00 & B & CA & 60000.00 & rent \\ 
  3 & 25000 & 26.30 & 36.00 & E & SC & 75000.00 & mortgage \\
  %1 & 7500 & 7.34 & 36 & A & MD & 70000 & rent \\
  %2 & 25000 & 9.43 & 60 & B & OH & 254000 & mortgage \\
  %3 & 14500 & 6.08 & 36 & A & MO & 80000 & mortgage \\
  $\vdots$ & $\vdots$ & $\vdots$ & $\vdots$ & $\vdots$ & $\vdots$
      & $\vdots$ & $\vdots$ \\
  50 & 15000 & 6.08 & 36.00 & A & TX & 77500.00 & mortgage \\
  %50 & 3000 & 7.96 & 36 & A & CA & 34000 & rent \\
   \hline
\end{tabular}
}
\caption{Empat baris dari matriks data \data{loan50}.}
\label{loan50DF}
\end{figure}
% Dropped: state, verified_income
% library(openintro); vars <- c("loan_amount", "interest_rate", "term", "grade", "state", "total_income", "homeownership"); library(xtable); data(loan50); loan50[c(1,2,3,50), vars]; xtable(loan50[c(1,2,3,50), vars])

\begin{figure}[h]
\centering\small
\begin{tabular}{lp{10.5cm}}
\hline
 {\bf variabel} & {\bf deskripsi} \\
\hline
\var{loan\us{}amount} & Jumlah pinjaman yang diterima,
    dalam dolar AS.  \\
\var{interest\us{}rate} & Tingkat bunga pinjaman,
    dalam persentase tahunan.  \\
\var{term} & Jangka waktu pinjaman, yang selalu dinyatakan
    sebagai bilangan bulat bulan. \\
\var{grade} & Peringkat pinjaman, yang bernilai A hingga G
    dan mencerminkan kualitas pinjaman serta kemungkinan
    pelunasannya.  \\
\var{state} & Negara bagian AS tempat peminjam berdomisili. \\
\var{total\us{}income} & Jumlah pendapatan peminjam,
    termasuk pendapatan tambahan apa pun, dalam dolar AS.   \\
\var{homeownership} & Menunjukkan apakah orang tersebut
    memiliki rumah tanpa hipotek, memiliki rumah dengan hipotek, atau menyewa.  \\
%\var{verified\us{}income} & Indicates whether the
%    income is verified, its source is verified but not the amount,
%    or it is not verified.   \\
\hline
\end{tabular}
\caption{Variabel dan deskripsinya dalam himpunan data \data{loan50}.}
\label{loan50Variables}
\end{figure}

\index{data!loan50|)}

Data pada Gambar~\ref{loan50DF} membentuk sebuah \term{matriks data},
yakni cara yang praktis dan lazim untuk mengorganisasi data,
terutama jika data dikumpulkan dalam lembar kerja.
Setiap baris matriks data bersesuaian dengan satu kasus
(unit observasi),
dan setiap kolom bersesuaian dengan satu variabel.
%A data matrix for the stroke study introduced in
%Section~\ref{basicExampleOfStentsAndStrokes} is shown
%in Figure~\vref{stentStudyResultsDF}, where the cases were
%patients and three variables were recorded for each
%patient.

Ketika mencatat data, gunakanlah matriks data kecuali terdapat
alasan yang sangat kuat untuk memakai struktur lain.
Struktur ini memungkinkan kasus baru ditambahkan sebagai baris
atau variabel baru sebagai kolom.

\begin{exercisewrap}
\begin{nexercise}
Nilai tugas, kuis, dan ujian dalam suatu mata kuliah sering
dicatat dalam buku nilai yang berbentuk matriks data.
Bagaimana Anda dapat mengorganisasi data nilai dengan menggunakan
matriks data?\footnotemark
\end{nexercise}
\end{exercisewrap}
\footnotetext{Ada beberapa strategi yang dapat digunakan.
  Salah satu strategi yang lazim adalah mewakili setiap mahasiswa dengan satu baris,
  lalu menambahkan satu kolom untuk setiap tugas, kuis, atau ujian.
  Dengan susunan ini, satu baris dapat ditinjau dengan mudah untuk memahami
  riwayat nilai seorang mahasiswa.
  Perlu pula disediakan kolom-kolom untuk informasi mahasiswa,
  misalnya satu kolom yang memuat nama mahasiswa.}

\index{data!county|(}

\begin{exercisewrap}
\begin{nexercise}\label{desc_county_as_data_matrix}%
Kita meninjau data untuk 3,142 county di Amerika Serikat,
yang mencakup nama setiap county,
negara bagiannya, populasinya pada 2017,
perubahan populasinya dari 2010 hingga 2017,
tingkat kemiskinan,
dan enam karakteristik tambahan.
Bagaimana data ini dapat diorganisasi dalam
sebuah matriks data?\footnotemark
\end{nexercise}
\end{exercisewrap}
\footnotetext{Setiap county dapat dipandang sebagai satu kasus,
  dan terdapat sebelas butir informasi yang dicatat untuk
  setiap kasus.
  Data ini dapat ditampung dalam tabel dengan 3,142 baris dan 11 kolom;
  setiap baris mewakili satu county dan setiap kolom mewakili
  satu butir informasi tertentu.}

Data yang dijelaskan dalam Latihan
Terbimbing~\ref{desc_county_as_data_matrix} membentuk himpunan data
\data{county}, yang ditampilkan sebagai matriks data
pada Gambar~\ref{countyDF}.
Variabel-variabelnya dirangkum pada Gambar~\ref{countyVariables}.

\begin{landscape}
\begin{figure}
\centering\small
\begin{tabular}{ccc ccc ccc ccc}
  \hline
 & \var{name} & \var{state} & \var{pop} & \var{pop\us{}change} & \var{poverty} & \var{homeownership} & \var{multi\us{}unit} & \var{unemp\us{}rate} & \var{metro} & \var{median\us{}edu} & \var{median\us{}hh\us{}income} \\ 
  \hline
1 & Autauga  & Alabama &  55504 &  1.48 & 13.7 & 77.5 &  7.2 & 3.86 & yes & some\us{}college & 55317 \\ 
  2 & Baldwin  & Alabama & 212628 &  9.19 & 11.8 & 76.7 & 22.6 & 3.99 & yes & some\us{}college & 52562 \\ 
  3 & Barbour  & Alabama &  25270 & -6.22 & 27.2 & 68.0 & 11.1 & 5.90 & no  & hs\us{}diploma   & 33368 \\ 
  4 & Bibb     & Alabama &  22668 &  0.73 & 15.2 & 82.9 &  6.6 & 4.39 & yes & hs\us{}diploma   & 43404 \\ 
  5 & Blount   & Alabama &  58013 &  0.68 & 15.6 & 82.0 &  3.7 & 4.02 & yes & hs\us{}diploma   & 47412 \\ 
  6 & Bullock  & Alabama &  10309 & -2.28 & 28.5 & 76.9 &  9.9 & 4.93 & no  & hs\us{}diploma   & 29655 \\ 
  7 & Butler   & Alabama &  19825 & -2.69 & 24.4 & 69.0 & 13.7 & 5.49 & no  & hs\us{}diploma   & 36326 \\ 
  8 & Calhoun  & Alabama & 114728 & -1.51 & 18.6 & 70.7 & 14.3 & 4.93 & yes & some\us{}college & 43686 \\ 
  9 & Chambers & Alabama &  33713 & -1.20 & 18.8 & 71.4 &  8.7 & 4.08 & no  & hs\us{}diploma   & 37342 \\ 
  10 & Cherokee & Alabama &  25857 & -0.60 & 16.1 & 77.5 &  4.3 & 4.05 & no  & hs\us{}diploma   & 40041 \\ 
  $\vdots$ & $\vdots$ & $\vdots$ & $\vdots$ & $\vdots$ & $\vdots$ & $\vdots$ & $\vdots$ & $\vdots$ & $\vdots$ & $\vdots$ & $\vdots$ \\
  3142 & Weston   & Wyoming &   6927 & -2.93 & 14.4 & 77.9 &  6.5 & 3.98 & no  & some\us{}college & 59605 \\ 
   \hline
\end{tabular}
\caption{Sebelas baris dari himpunan data \data{county}.}
\label{countyDF}
% library(openintro); data(county); county$name <- gsub(" County$", "", county$name); county$pop <- county$pop2017; county$unemp_rate = county$unemployment_rate; these <- c("name", "state", "pop", "pop_change", "poverty", "homeownership", "multi_unit", "unemp_rate", "metro", "median_edu", "median_hh_income"); county <- county[, these]; library(xtable); xtable(as.data.frame(lapply(rbind.data.frame(head(county, 10), tail(county, 1)), function(x) { format(x) })))
\end{figure}

\begin{figure}
\centering\small
\begin{tabular}{lp{11cm}}
\hline
{\bf variabel} & {\bf deskripsi} \\
\hline
\var{name} &
    Nama county. \\
\var{state} &
    Negara bagian tempat county berada,
    atau Distrik Columbia. \\
\var{pop} &
    Populasi pada 2017. \\
\var{pop\us{}change} &
    Persentase perubahan populasi dari 2010 hingga 2017.
    Sebagai contoh, nilai \resp{1.48} pada baris pertama
    berarti populasi county tersebut
    meningkat sebesar 1.48\% dari 2010 hingga 2017. \\
\var{poverty} &
    Persentase populasi yang hidup dalam kemiskinan. \\
\var{homeownership}  &
    Persentase populasi yang tinggal di rumah milik sendiri
    atau tinggal bersama pemiliknya, misalnya anak yang tinggal bersama orang tua
    pemilik rumah. \\
\var{multi\us{}unit}  &
    Persentase unit tempat tinggal yang berada dalam bangunan multiunit,
    misalnya apartemen. \\
\var{unemp\us{}rate} &
    Tingkat pengangguran dalam persen. \\
\var{metro} &
    Menunjukkan apakah county tersebut mencakup kawasan metropolitan. \\
\var{median\us{}edu} & Tingkat pendidikan median, yang
    dapat bernilai salah satu dari
    \resp{below\us{}hs},
    \resp{hs\us{}diploma},
    \resp{some\us{}college},
    dan \resp{bachelors}. \\
\var{median\us{}hh\us{}income} &
    Median pendapatan rumah tangga di county tersebut; pendapatan rumah tangga
    sama dengan jumlah pendapatan para penghuninya yang berusia
    15~tahun atau lebih. \\
%\var{per\us{}capita\us{}income} &
%    Per capita (per person) income for the county. \\
\hline
\end{tabular}
\centering
\caption{Variabel dan deskripsinya dalam himpunan data \data{county}.}
\label{countyVariables}
\end{figure}
\end{landscape}

\subsection{Jenis variabel}
\label{variableTypes}

Perhatikan variabel \var{unemp\us{}rate}, \var{pop}, \var{state},
dan \var{median\us{}edu} dalam himpunan data \data{county}.
Masing-masing variabel ini pada dasarnya berbeda dari ketiga variabel
lainnya, tetapi beberapa di antaranya memiliki karakteristik tertentu yang sama.

Pertama, perhatikan \var{unemp\us{}rate},
yang disebut variabel \term{numerik} karena dapat
mengambil beragam nilai numerik,
dan penjumlahan, pengurangan, atau pengambilan rata-rata
atas nilai-nilai tersebut masuk akal.
Sebaliknya, variabel yang mencatat kode area telepon tidak akan
kita klasifikasikan sebagai numerik karena rata-rata, jumlah, dan
selisih kode area tidak memiliki makna yang jelas.

Variabel \var{pop} juga bersifat numerik, meskipun tampak
sedikit berbeda dari \var{unemp\us{}rate}.
Variabel jumlah populasi ini hanya dapat mengambil bilangan bulat
nonnegatif (\resp{0}, \resp{1}, \resp{2}, ...).
Karena alasan ini, variabel populasi disebut \term{diskret}
sebab hanya dapat mengambil nilai numerik dengan lompatan tertentu.
Sebaliknya, variabel tingkat pengangguran disebut \term{kontinu}.

Variabel \var{state} dapat mengambil hingga 51 nilai jika
Washington, DC turut dihitung: \resp{AL}, \resp{AK}, ...,
dan \resp{WY}.
Karena responsnya sendiri berupa kategori,
\var{state} disebut variabel \term{kategoris},
dan nilai-nilai yang mungkin disebut \term{level} variabel tersebut.

Terakhir, perhatikan variabel \var{median\us{}edu},
yang menggambarkan tingkat pendidikan median penduduk county
dan mengambil nilai
\resp{below\us{}hs}, \resp{hs\us{}diploma},
\resp{some\us{}college}, atau \resp{bachelors}
di setiap county.
Variabel ini tampak sebagai gabungan: variabel tersebut bersifat kategoris,
tetapi level-levelnya memiliki urutan alami.
Variabel dengan sifat-sifat ini disebut variabel \term{ordinal},
sedangkan variabel kategoris biasa tanpa urutan khusus semacam ini
disebut variabel \term{nominal}.
Untuk menyederhanakan analisis, setiap variabel ordinal dalam buku ini
akan diperlakukan sebagai variabel kategoris nominal (tidak berurutan).

\begin{figure}[h]
  \centering
  \Figure
    [Pengelompokan variabel menurut jenisnya: "semua variabel" terbagi menjadi "numerik" dan "kategoris". Selanjutnya, "numerik" terbagi menjadi "kontinu" dan "diskret", sedangkan "kategoris" terbagi menjadi variabel "nominal (kategoris tidak berurutan)" dan "ordinal (kategoris berurutan)".]
    {0.57}{variables}
  \caption{Pengelompokan variabel menurut jenisnya.}
  \label{variables}
\end{figure}

\begin{examplewrap}
\begin{nexample}{Data dikumpulkan mengenai mahasiswa
    dalam suatu mata kuliah statistika.
    Tiga variabel dicatat untuk setiap mahasiswa:
    jumlah saudara kandung, tinggi badan mahasiswa, dan apakah
    mahasiswa tersebut pernah mengambil mata kuliah statistika.
    Klasifikasikan setiap variabel sebagai numerik kontinu,
    numerik diskret, atau kategoris.}
  Jumlah saudara kandung dan tinggi badan mahasiswa merupakan
  variabel numerik.
  Karena jumlah saudara kandung merupakan hitungan, variabel ini diskret.
  Tinggi badan bervariasi secara kontinu, sehingga merupakan variabel
  numerik kontinu.
  Variabel terakhir mengklasifikasikan mahasiswa ke dalam dua kategori
  -- yang pernah dan yang belum pernah mengambil mata kuliah statistika --
  sehingga variabel tersebut bersifat kategoris.
\end{nexample}
\end{examplewrap}

\begin{exercisewrap}
\begin{nexercise}\index{data!migrain}%
Sebuah eksperimen mengevaluasi efektivitas obat baru
untuk menangani migrain.
Variabel \var{group} digunakan untuk menunjukkan kelompok eksperimen
setiap pasien: perlakuan atau kontrol.
Variabel \mbox{\var{num\us{}migraines}} menyatakan jumlah
migrain yang dialami pasien selama periode 3 bulan.
\mbox{Klasifikasikan} setiap variabel sebagai numerik atau
kategoris.\footnotemark
\end{nexercise}
\end{exercisewrap}
\footnotetext{Variabel
  \var{group} hanya dapat mengambil salah satu dari dua nama kelompok,
  sehingga bersifat kategoris.
  Variabel \var{num\us{}migraines} menggambarkan
  hitungan jumlah migrain, yakni luaran yang dapat dikenai
  operasi aritmetika dasar secara bermakna; karena itu, luaran ini bersifat
  numerik. Lebih khusus lagi, karena menyatakan hitungan,
  \var{num\us{}migraines} merupakan variabel numerik diskret.}


\D{\newpage}

\subsection{Hubungan antarvariabel}
\label{variableRelations}

Banyak analisis berawal dari upaya peneliti mencari
hubungan antara dua variabel atau lebih.
Seorang ilmuwan sosial mungkin ingin menjawab beberapa
pertanyaan berikut:
\newcommand{\popchangevmedianhhincomequestion}[0]{
    % Note that this question is used to introduce the
    %explanatory / response variable topic.
    Apakah kenaikan populasi county yang melebihi rata-rata
    cenderung bersesuaian dengan county yang memiliki median
    pendapatan rumah tangga lebih tinggi atau lebih rendah?}%
\begin{enumerate}
\setlength{\itemsep}{0mm}
\item[(1)]\label{ownershipMultiUnitQuestion}
    Jika tingkat kepemilikan rumah di suatu county lebih rendah
    daripada rata-rata nasional, apakah persentase bangunan multiunit
    di county tersebut cenderung berada di atas atau di bawah rata-rata nasional?
\item[(2)]\label{pop_change_v_median_hh_income_question}
    \popchangevmedianhhincomequestion{}
    % Do counties with a higher median household income
    % tend to be growing faster or slower than other counties?
\item[(3)]\label{isAverageIncomeAssociatedWithSmokingBans}
    Seberapa bergunakah tingkat pendidikan median sebagai prediktor
    median pendapatan rumah tangga county di AS?
\end{enumerate}

Untuk menjawab pertanyaan-pertanyaan ini, data harus dikumpulkan, misalnya
himpunan data \data{county} yang ditampilkan pada Gambar~\ref{countyDF}.
Pemeriksaan statistik ringkasan \index{statistik ringkasan} dapat
memberikan wawasan bagi ketiga pertanyaan tentang county tersebut.
Selain itu, grafik dapat digunakan untuk menjelajahi data secara visual.

\indexthis{Diagram pencar}{diagram pencar} merupakan salah satu jenis grafik
yang digunakan untuk mempelajari hubungan antara dua variabel numerik.
Gambar~\ref{multiunitsVsOwnership} membandingkan variabel
\var{homeownership} dan
\var{multi\us{}unit},
yang merupakan persentase unit dalam bangunan multiunit
(misalnya apartemen dan kondominium).
Setiap titik pada diagram mewakili satu county.
Sebagai contoh, titik yang disorot bersesuaian dengan
County~413 dalam himpunan data \data{county}:
Chattahoochee County, Georgia, dengan 39.4\% unit
berada dalam bangunan multiunit dan tingkat kepemilikan rumah
sebesar 31.3\%.
Diagram pencar tersebut menunjukkan hubungan antara kedua
variabel: county dengan tingkat bangunan multiunit yang lebih tinggi
cenderung memiliki tingkat kepemilikan rumah yang lebih rendah.
Kita dapat memikirkan berbagai kemungkinan penyebab hubungan ini,
lalu menyelidiki setiap gagasan untuk menentukan penjelasan yang
paling masuk akal.

\begin{figure}[h]
  \centering
  \Figure
    [Diagram pencar ribuan county dengan persentase bangunan multiunit di setiap county pada sumbu horizontal dan tingkat kepemilikan rumah pada sumbu vertikal. Data kedua variabel berkisar dari 0\% hingga hampir 100\%. Secara umum, titik-titik jauh lebih terkonsentrasi di sudut kiri atas grafik, lalu cenderung menurun dan semakin jarang untuk observasi yang lebih jauh ke kanan. Satu titik diberi anotasi pada lokasi (39.4\%, 31.3\%).]
    {0.79}{multiunitsVsOwnership}
  \caption{Diagram pencar tingkat kepemilikan rumah terhadap persentase
      unit dalam bangunan multiunit untuk county di AS.
      Titik yang disorot mewakili Chattahoochee County, Georgia,
      yang memiliki tingkat bangunan multiunit 39.4\% dan tingkat kepemilikan rumah
      31.3\%.}
  \label{multiunitsVsOwnership}
\end{figure}

Tingkat bangunan multiunit dan kepemilikan rumah dikatakan
berasosiasi karena diagram menunjukkan pola yang dapat dikenali.
Jika dua variabel menunjukkan suatu keterkaitan satu sama lain,
keduanya disebut variabel yang \term{berasosiasi}.
Variabel yang berasosiasi juga dapat disebut variabel
\termsub{terikat}{variabel terikat},
dan sebaliknya.

\D{\newpage}

\begin{exercisewrap}
\begin{nexercise}
Perhatikan variabel-variabel dalam himpunan data \data{loan50},
yang dijelaskan pada Gambar~\vref{loan50Variables}.
Buatlah dua pertanyaan yang menarik bagi Anda tentang kemungkinan
hubungan antarvariabel dalam \data{loan50}.\footnotemark
\end{nexercise}
\end{exercisewrap}
\footnotetext{Dua contoh pertanyaan:
  (1)~Bagaimana hubungan antara jumlah pinjaman dan
      pendapatan total?
  (2)~Jika pendapatan seseorang berada di atas rata-rata, apakah
      tingkat bunganya cenderung berada di atas atau di bawah rata-rata?}

\begin{examplewrap}
\begin{nexample}{Contoh ini menelaah hubungan
    antara perubahan populasi suatu county
    dari 2010 hingga 2017
    dan median pendapatan rumah tangga,
    yang divisualisasikan sebagai diagram pencar pada
    Gambar~\ref{pop_change_v_med_income}.
    Apakah kedua variabel ini berasosiasi?}
  Semakin besar median pendapatan rumah tangga suatu county,
  semakin tinggi pertumbuhan populasi yang diamati di county tersebut.
  Meskipun tren ini tidak berlaku bagi setiap county,
  tren pada diagram tampak jelas.
  Karena terdapat suatu hubungan antara kedua variabel,
  keduanya berasosiasi.
\end{nexample}
\end{examplewrap}

%When two variables show some connection with one another,
%they are called \term{associated} variables.
%Associated variables can also be called \term{dependent}
%variables and vice-versa.
%When the variables increase together,
%as they do in Figure~\ref{loan_amount_vs_income},
%they are said to be \term{positively associated}.
%When the trend in the scatterplot goes down to the right,
%then they are described as \term{negatively correlated}.

%While we may find it interesting to consider the relationship
%between two variables such as those in the scatterplot,
%the relationship between those variables can be more complex.
%For example, interest rates on loans tend to be chosen based
%on the riskiness of the loan, i.e. how likely it is to be
%paid back, and that is likely to depend on a variety of
%details, such as what the loan is for, the person's
%creditworthiness, whether their income is verified, etc.
%We will begin exploring some of these more complex relationships
%in graphs in Chapter~\ref{ch_summarizing_data} and beyond.
%\Comment{Revise if we don't add these more rich plots...}

%\begin{example}{Figure~\ref{interest_rate_vs_loan_amount}
%    features a scatterplot of interest rate against loan amount.
%    Are these variables associated?}
%  There isn't an evident trend in the data,
%  so we would say these two variables are not associated.
%\end{example}

\begin{figure}
  \centering
  \Figure
    [Diagram pencar ribuan county dengan median pendapatan rumah tangga pada sumbu horizontal (data berkisar dari \$0 hingga \$120,000) dan perubahan populasi selama 7 tahun pada sumbu vertikal (data berkisar dari sekitar -15\% hingga 25\%). Terdapat gugus titik yang berpusat di sekitar (\$45,000, -1\%); titik-titik menunjukkan sedikit tren naik sekaligus menjadi semakin jarang dan lebih berfluktuasi pada observasi dengan median pendapatan lebih tinggi. Satu titik diberi anotasi pada lokasi (\$22,736, -3.63\%).]
    {0.9}{pop_change_v_med_income}
  \caption{Diagram pencar yang menampilkan
      \var{pop\us{}change}
      terhadap \var{median\us{}hh\us{}income}.
      Owsley County di Kentucky disorot;
      county tersebut kehilangan 3.63\% populasinya dari 2010 hingga 2017
      dan memiliki median pendapatan rumah tangga sebesar \$22,736.}
  \label{pop_change_v_med_income}
\end{figure}

Karena terdapat tren menurun pada
Gambar~\ref{multiunitsVsOwnership} --
county dengan lebih banyak unit dalam bangunan multiunit
berasosiasi dengan kepemilikan rumah yang lebih rendah --
kedua variabel ini dikatakan memiliki
\termsub{asosiasi negatif}{asosiasi negatif}.
Sebuah \term{asosiasi positif} tampak dalam hubungan antara
\var{median\us{}hh\us{}income}
dan \var{pop\us{}change}
pada Gambar~\ref{pop_change_v_med_income},
di mana county dengan median pendapatan rumah tangga yang lebih tinggi
cenderung memiliki tingkat pertumbuhan populasi yang lebih tinggi.

Jika dua variabel tidak berasosiasi,
keduanya dikatakan \termsub{bebas}{variabel bebas} satu sama lain.
Dengan kata lain, dua variabel saling bebas jika tidak terdapat
hubungan yang jelas di antara keduanya.

\begin{onebox}{Berasosiasi atau bebas, bukan keduanya}
Sepasang variabel memiliki suatu keterkaitan (berasosiasi) atau tidak memiliki keterkaitan (saling bebas). Tidak ada pasangan variabel yang sekaligus berasosiasi dan saling bebas.
\end{onebox}


\D{\newpage}

\subsection{Variabel penjelas dan variabel respons}
\label{explanatoryAndResponse}

Ketika mengajukan pertanyaan tentang hubungan antara dua variabel,
kadang-kadang kita juga ingin menentukan apakah perubahan pada satu
variabel menyebabkan perubahan pada variabel lainnya.
Perhatikan perumusan ulang berikut dari pertanyaan sebelumnya
tentang himpunan data \data{county}:
\begin{quote}\em
  Jika median pendapatan rumah tangga di suatu county meningkat,
  apakah hal ini mendorong peningkatan populasinya?
\end{quote}
Dalam pertanyaan ini, kita menanyakan apakah satu variabel
memengaruhi variabel lainnya.
Jika itulah dugaan yang mendasari pertanyaan kita,
maka \emph{median pendapatan rumah tangga} merupakan
\termsub{variabel penjelas}{variabel penjelas}, sedangkan
\emph{perubahan populasi} merupakan
\termsub{variabel respons}{variabel respons}
dalam hubungan yang dihipotesiskan.\footnote{Variabel penjelas dan respons
  kadang disebut variabel \termsub{bebas}{variabel bebas} dan
  \termsub{terikat}{variabel terikat}. Karena istilah bebas/terikat juga dapat
  menyatakan hubungan dalam suatu \emph{pasangan} variabel, buku ini
  menghindari sebutan tersebut.}

\index{data!county|)}

\begin{onebox}{Variabel penjelas dan variabel respons}
Ketika kita menduga bahwa satu variabel mungkin memengaruhi variabel lain
secara kausal, variabel pertama disebut variabel penjelas dan variabel
kedua disebut variabel respons.
\vspace{1mm}

\hspace{10mm}\Figure
    [Grafik sederhana yang menampilkan kata "variabel penjelas" dengan panah menuju "variabel respons"; kata "mungkin memengaruhi" berada di atas panah.]
    {0.53}{expResp}

Bagi banyak pasangan variabel, tidak ada hubungan yang dihipotesiskan;
dalam keadaan seperti itu, kedua label ini tidak diterapkan pada
variabel mana pun.
\end{onebox}

Pemberian label ini tidak menjamin adanya hubungan kausal.
Dalam pengantar ini, eksperimen merupakan cara paling langsung untuk memeriksa
efek kausal. Studi nonacak juga dapat mendukung inferensi kausal, tetapi hanya
dengan asumsi dan metode tambahan.

%\begin{exercisewrap}
%\begin{nexercise}
%Consider the earlier question:
%\begin{quote}\em
%  If a county has a higher median household income,
%  does this drive an increase in its population?
%\end{quote}
%We could have just as easily reframed the causal relationship
%to be in the reverse direction:
%\begin{quote}\em
%  If a county more population growth, does this drive
%  it to have a higher median household income?
%\end{quote}
%What are the explanatory and response variables when framing
%the variable relationship in the second question?
%\end{nexercise}
%\end{exercisewrap}
%\footnotetext{In this framing, we have hypothesized
%  that population growth drives median household income.
%  That is, population growth is the explanatory variable,
%  and median household income is the response.
%  This exercise should emphasize that these variable labels
%  do not actually define whether one variable actually affects
%  the other.}


\subsection{Pengantar studi observasional dan eksperimen}

\noindent%
Terdapat dua jenis utama pengumpulan data:
studi observasional dan eksperimen.

Dalam \term{studi observasional}, peneliti mengumpulkan data tanpa campur
tangan langsung dalam proses pembentukannya---misalnya melalui survei,
rekam medis atau catatan perusahaan, atau dengan mengikuti sebuah
\term{kohort} individu serupa untuk menyusun hipotesis tentang penyebab
penyakit. Peneliti hanya mengamati data yang timbul. Studi seperti ini dapat
menunjukkan asosiasi alami antarvariabel, tetapi tidak dengan sendirinya
membuktikan hubungan kausal.

Untuk menyelidiki kemungkinan hubungan kausal, peneliti dapat melakukan
\term{eksperimen}: sebuah sampel dibagi ke dalam kelompok-kelompok yang
masing-masing \emph{diberi} perlakuan. Biasanya terdapat variabel penjelas dan
variabel respons. Misalnya, kita dapat menguji dugaan bahwa obat menurunkan
angka kematian pasien serangan jantung selama setahun. Jika individu
dialokasikan secara acak ke dalam kelompok, desain tersebut disebut
\term{eksperimen acak}. Dalam uji obat tadi, pelemparan koin dapat menentukan
apakah pasien menerima \term{plasebo} (perlakuan semu) atau obat.
Lihat studi kasus pada
Bagian~\ref{basicExampleOfStentsAndStrokes} untuk contoh eksperimen
lain, meskipun penelitian tersebut tidak menggunakan plasebo.

\begin{onebox}{Asosiasi $\neq$ Kausalitas}
  Secara umum, asosiasi tidak dengan sendirinya menyiratkan kausalitas.
  Eksperimen acak memberikan dasar paling langsung untuk inferensi kausal
  dalam pembahasan ini; desain lain memerlukan asumsi tambahan yang eksplisit.
\end{onebox}


{\exercisesheader{}

% 3

\eoce{\qt{Polusi udara dan luaran kelahiran, komponen penelitian\label{study_components_airpoll}}
Para peneliti mengumpulkan data untuk mengkaji hubungan antara polutan udara
dan kelahiran prematur di California Selatan. Selama penelitian, tingkat polusi
udara diukur oleh stasiun pemantau kualitas udara. Secara khusus, kadar karbon
monoksida dicatat dalam bagian per juta, nitrogen dioksida dan ozon dalam bagian
per seratus juta, dan partikulat kasar (PM$_{10}$) dalam $\mu g/m^3$.
Data lama kehamilan dikumpulkan untuk 143,196 kelahiran antara tahun 1989
dan 1993, lalu paparan polusi udara selama kehamilan dihitung untuk setiap
kelahiran. Analisis menunjukkan bahwa peningkatan PM$_{10}$ di udara sekitar dan,
pada tingkat yang lebih rendah, konsentrasi CO mungkin berkaitan dengan kejadian
kelahiran prematur.\footfullcite{Ritz+Yu+Chapa+Fruin:2000}
\begin{parts}
\item Tentukan pertanyaan penelitian utama dalam penelitian ini.
\item Siapa subjek dalam penelitian ini, dan berapa banyak yang disertakan?
\item Apa saja variabel dalam penelitian ini? Tentukan apakah setiap variabel
bersifat numerik atau kategoris. Jika numerik, nyatakan apakah variabel tersebut
diskret atau kontinu. Jika kategoris, nyatakan apakah variabel tersebut ordinal.
\end{parts}
}{}

% 4

\eoce{\qt{Metode Buteyko, komponen penelitian\label{study_components_buteyko}}
Metode Buteyko adalah teknik pernapasan dangkal yang dikembangkan oleh Konstantin
Buteyko, seorang dokter Rusia, pada tahun 1952. Bukti anekdotal menunjukkan bahwa
metode Buteyko dapat mengurangi gejala asma dan meningkatkan kualitas hidup. Dalam
sebuah penelitian ilmiah untuk menentukan keefektifan metode ini, para peneliti
merekrut 600 pasien asma berusia 18-69 tahun yang bergantung pada obat untuk
penanganan asma. Pasien-pasien ini secara acak dibagi menjadi dua kelompok
penelitian: satu kelompok mempraktikkan metode Buteyko dan kelompok lainnya tidak.
Pasien diberi skor dari 0 hingga 10 untuk kualitas hidup, aktivitas, gejala asma,
dan pengurangan penggunaan obat. Rata-rata, peserta dalam kelompok Buteyko
mengalami penurunan gejala asma yang signifikan dan peningkatan kualitas hidup.\footfullcite{McDowan:2003}
\begin{parts}
\item Tentukan pertanyaan penelitian utama dalam penelitian ini.
\item Siapa subjek dalam penelitian ini, dan berapa banyak yang disertakan?
\item Apa saja variabel dalam penelitian ini? Tentukan apakah setiap variabel
bersifat numerik atau kategoris. Jika numerik, nyatakan apakah variabel tersebut
diskret atau kontinu. Jika kategoris, nyatakan apakah variabel tersebut ordinal.
\end{parts}
}{}

% 5

\eoce{\qt{Kecurangan, komponen penelitian\label{study_components_cheaters}}
Para peneliti yang mengkaji hubungan antara kejujuran, usia, dan pengendalian diri
melakukan eksperimen terhadap 160 anak berusia antara 5 dan 15 tahun.
Peserta melaporkan usia, jenis kelamin, dan apakah mereka anak tunggal.
Para peneliti meminta setiap anak untuk melempar koin seimbang secara pribadi,
mencatat hasilnya (putih atau hitam) pada selembar kertas, dan memberi tahu bahwa
hanya anak yang melaporkan putih yang akan mendapat hadiah.
Temuan penelitian dapat dirangkum sebagai berikut:
``Separuh anak secara tegas diberi tahu agar tidak berbuat curang, sedangkan
yang lain tidak diberi petunjuk tegas apa pun.
Dalam kelompok tanpa petunjuk, probabilitas berbuat curang ditemukan seragam
di antara kelompok yang dibentuk berdasarkan karakteristik anak.
Dalam kelompok yang secara tegas diberi tahu agar tidak berbuat curang,
anak perempuan lebih kecil kemungkinannya berbuat curang; sementara tingkat
kecurangan tidak bervariasi menurut usia pada anak laki-laki, tingkat tersebut
menurun seiring bertambahnya usia pada anak perempuan.''\footfullcite{Bucciol:2011}
\begin{parts}
\item Tentukan pertanyaan penelitian utama dalam penelitian ini.
\item Siapa subjek dalam penelitian ini, dan berapa banyak yang disertakan?
\item
  Berapa banyak variabel yang dicatat untuk setiap subjek dalam penelitian ini
  agar temuan tersebut dapat disimpulkan? Nyatakan variabel-variabel itu beserta
  jenisnya.
\end{parts}
}{}

\D{\newpage}

% 6

\eoce{\qt{Perilaku mengambil permen, komponen penelitian\label{study_components_stealers}}
Dalam sebuah penelitian tentang hubungan antara kelas sosioekonomi dan perilaku
tidak etis, 129 mahasiswa sarjana University of California, Berkeley diminta
mengidentifikasi diri sebagai anggota kelas sosial bawah atau atas dengan
membandingkan diri mereka dengan orang lain yang memiliki uang paling banyak
(paling sedikit), pendidikan paling tinggi (paling rendah), dan pekerjaan paling
dihormati (paling tidak dihormati). Mereka juga diberi sebuah stoples berisi
permen yang dibungkus satu per satu dan diberi tahu bahwa permen tersebut
diperuntukkan bagi anak-anak di laboratorium terdekat, tetapi mereka boleh
mengambil beberapa jika mau. Setelah menyelesaikan beberapa tugas yang tidak
berkaitan, peserta melaporkan jumlah permen yang mereka ambil.\footfullcite{Piff:2012}
\begin{parts}
\item Tentukan pertanyaan penelitian utama dalam penelitian ini.
\item Siapa subjek dalam penelitian ini, dan berapa banyak yang disertakan?
\item Penelitian menemukan bahwa mahasiswa yang diidentifikasi sebagai kelas atas
mengambil lebih banyak permen daripada mahasiswa lainnya. Berapa banyak variabel
yang dicatat untuk setiap subjek dalam penelitian ini agar temuan tersebut dapat
disimpulkan? Nyatakan variabel-variabel itu beserta jenisnya.
\end{parts}
}{}

% 7

\eoce{\qt{Migrain dan akupunktur,
    Bagian II\label{migraine_and_acupuncture_exp_resp}}
Latihan~\ref{migraine_and_acupuncture_intro}
memperkenalkan sebuah penelitian yang mengkaji apakah akupunktur
memiliki pengaruh terhadap migrain.
Para peneliti melakukan penelitian terkontrol acak
dengan menempatkan pasien secara acak ke salah satu dari dua kelompok:
perlakuan atau kontrol.
Pasien dalam kelompok perlakuan menerima akupunktur
yang dirancang secara khusus untuk menangani migrain.
Pasien dalam kelompok kontrol menerima akupunktur plasebo
(jarum ditusukkan pada lokasi yang bukan titik akupunktur).
Pasien ditanya 24 jam setelah menerima akupunktur
apakah mereka sudah bebas nyeri.
Apa variabel penjelas dan variabel respons dalam penelitian ini?
}{}

% 8

\eoce{\qt{Sinusitis dan antibiotik,
    Bagian II\label{sinusitis_and_antibiotics_exp_resp}}
Latihan~\ref{sinusitis_and_antibiotics_intro}
memperkenalkan sebuah penelitian yang mengkaji pengaruh pengobatan antibiotik
terhadap sinusitis akut.
Peserta penelitian menerima rangkaian pengobatan antibiotik selama 10 hari
(perlakuan) atau plasebo yang tampak dan terasa serupa (kontrol).
Pada akhir periode 10 hari, pasien ditanya apakah
gejala mereka membaik.
Apa variabel penjelas dan variabel respons dalam penelitian ini?
}{}

% 9

\eoce{\qt{Bunga iris Fisher\label{fisher_irises}} Sir Ronald Aylmer Fisher adalah
seorang ahli statistika, biologi evolusi, dan genetika asal Inggris yang bekerja
dengan sebuah himpunan data berisi panjang dan lebar sepal serta panjang dan lebar
petal dari tiga spesies bunga iris (\textit{setosa}, \textit{versicolor}, dan
\textit{virginica}). Himpunan data tersebut memuat 50 bunga dari setiap spesies.
\footfullcite{Fisher:1936} \\
\noindent\begin{minipage}[c]{0.48\textwidth}
\begin{parts}
\item Berapa banyak kasus yang disertakan dalam data?
\item Berapa banyak variabel numerik yang disertakan dalam data? Sebutkan
variabel-variabel tersebut dan tentukan apakah masing-masing kontinu atau diskret.
\item Berapa banyak variabel kategoris yang disertakan dalam data, dan apa saja
variabel tersebut? Daftarkan level (kategori) yang bersesuaian.
\end{parts}
\end{minipage}
\begin{minipage}[c]{0.01\textwidth}
\ 
\end{minipage}
\begin{minipage}[c]{0.2\textwidth}
\begin{center}
\Figures[Foto bunga iris ungu.]{}{eoce/fisher_irises}{irisversicolor}
\end{center}
\end{minipage}
\begin{minipage}[c]{0.01\textwidth}
\ 
\end{minipage}
\begin{minipage}[c]{0.23\textwidth}
{\raggedright\footnotesize Foto oleh Ryan Claussen
(\oiRedirect{textbook-flickr_ryan_claussen_iris_picture}{http://flic.kr/p/6QTcuX}) 
\href{https://creativecommons.org/licenses/by-sa/2.0/}{lisensi CC~BY-SA~2.0}}
\end{minipage}
}{}

% 10

\eoce{\qt{Kebiasaan merokok penduduk Britania Raya\label{smoking_habits_UK_datamatrix}} Sebuah survei
dilakukan untuk mempelajari kebiasaan merokok penduduk Britania Raya. Matriks
data di bawah menampilkan sebagian data yang dikumpulkan dalam survei ini.
Perhatikan bahwa ``$\pounds$'' berarti pound sterling Britania, ``cig'' berarti
rokok, dan ``N/A'' menunjukkan komponen data yang hilang. \footfullcite{data:smoking}
\begin{center}
\scriptsize{
\begin{tabular}{rccccccc}
\hline
	& kelamin 	 & usia 	& status 	& pendapatan 					     & merokok & akhirPekan	& hariKerja \\
\hline
1 	& Female & 42 	& Single 	& Under $\pounds$2,600 			     & Yes 	 & 12 cig/day   & 12 cig/day \\ 
2 	& Male	 & 44	& Single 	& $\pounds$10,400 to $\pounds$15,600 & No	 & N/A 			& N/A \\ 
3 	& Male 	 & 53 	& Married   & Above $\pounds$36,400 		     & Yes 	 & 6 cig/day 	& 6 cig/day \\ 
\vdots & \vdots & \vdots & \vdots & \vdots 				             & \vdots & \vdots 	    & \vdots \\ 
1691 & Male  & 40   & Single 	& $\pounds$2,600 to $\pounds$5,200   & Yes 	 & 8 cig/day 	& 8 cig/day \\   
\hline
\end{tabular}
}
\end{center}
\begin{parts}
\item Apa yang diwakili oleh setiap baris matriks data?
\item Berapa banyak peserta yang disertakan dalam survei?
\item Tentukan apakah setiap variabel dalam penelitian bersifat numerik atau
kategoris. Jika numerik, tentukan apakah variabel itu kontinu atau diskret.
Jika kategoris, nyatakan apakah variabel tersebut ordinal.
\end{parts}
}{}

% 11

\eoce{\qt{Bandara di AS\label{US Airports}}
Visualisasi di bawah menunjukkan sebaran geografis bandara di Amerika Serikat
daratan utama dan Washington, DC.
Visualisasi ini dibuat berdasarkan sebuah himpunan data dengan
setiap amatan berupa satu bandara.\footfullcite{data:usairports}
Batas kartografis pada peta bersumber dari U.S. Census Bureau.
\par\noindent\makebox[\textwidth][c]{%
\Figures[Empat salinan peta Amerika Serikat ditampilkan dalam kisi 2 $\times$ 2. Pada setiap peta, sumbu diberi label bujur (130 derajat barat hingga 60 derajat barat) dan lintang (20 derajat utara hingga 50 derajat utara). Kolom pertama grafik diberi label "penggunaan privat" dan kolom kedua "penggunaan publik". Baris pertama grafik diberi label "milik privat" dan baris kedua "milik publik". Keempat grafik menampilkan titik-titik yang masing-masing mewakili satu bandara. Ribuan titik tampak pada peta kiri atas (penggunaan privat, milik privat) dan peta kanan bawah (penggunaan publik, milik publik), sedangkan dua peta lainnya menampilkan relatif lebih sedikit titik, meskipun jumlahnya masih berkisar dari ratusan hingga beberapa ribu. Pada semua grafik, titik lebih padat di bagian tengah dan timur Amerika Serikat, lebih jarang di wilayah pegunungan dan gurun, lalu kembali lebih padat di sekitar negara bagian yang berbatasan dengan Samudra Pasifik, terutama di dekat kota-kota besar.]{0.85}{eoce/airports}{airports}
}\par
\begin{parts}
\item
    Daftarkan variabel yang digunakan untuk membuat visualisasi ini.
\item
    Tentukan apakah setiap variabel dalam penelitian bersifat numerik
    atau kategoris.
    Jika numerik, tentukan apakah variabel itu kontinu atau diskret.
    Jika kategoris, nyatakan apakah variabel tersebut ordinal.
\end{parts}
}{}

% 12

\eoce{\qt{Pemungutan suara PBB\label{unvotes}}
Visualisasi berikut menunjukkan pola suara Amerika Serikat, Kanada, dan Meksiko
mengenai beragam isu di Majelis Umum Perserikatan Bangsa-Bangsa. Untuk setiap
tahun 1946--2015, visualisasi ini menampilkan persentase pemungutan suara
tercatat ketika setiap negara memilih ya pada setiap isu. Data terdiri atas
amatan pasangan negara--tahun.\footfullcite{data:unvotes}
\par\noindent\makebox[\textwidth][c]{%
\Figures[Ditampilkan sebuah kisi diagram pencar dengan garis tren yang ditumpangkan untuk masing-masing dari tiga kelompok titik (berwarna hijau, biru, dan merah) pada setiap grafik. Kisi grafik terdiri atas 2 baris dan 3 kolom; grafik dalam deskripsi ini disebut dengan nomor 1 sampai 3 pada baris pertama dan 4 sampai 6 pada baris kedua. Pada semua grafik, sumbu horizontal menyatakan "tahun" (sekitar 1945 hingga sekitar 2018), sedangkan sumbu vertikal menyatakan "persentase ya" dengan nilai dari 0\% hingga 100\%. Setiap grafik merangkum pola pemungutan suara untuk satu topik berbeda di Majelis Umum PBB dan untuk negara Kanada (biru), Meksiko (hijau), serta Amerika Serikat (merah). Setiap grafik memuat titik dan garis tren lentur (nonlinier) yang disesuaikan dengan titik-titik tersebut. Dalam semua kasus kecuali Grafik 2 untuk "Kolonialisme", titik (data) relatif jarang pada 1940 hingga 1960 dibandingkan tahun-tahun berikutnya. Grafik 1 mewakili "Pengendalian senjata dan perlucutan senjata"; semua negara bermula rendah, antara 0\% dan 25\%, lalu naik cepat menjelang 1960 menjadi antara 25\% dan 95\%. AS tetap paling rendah (berkisar sekitar 25\% hingga 40\%), Kanada sedikit lebih tinggi pada 50\% hingga 70\%, dan Meksiko paling tinggi, biasanya antara 85\% dan 100\%. Grafik 2 berlabel "Kolonialisme"; garis tren bermula antara 50\% dan 80\%, kemudian AS turun mendekati 0\% pada 1980, Kanada berfluktuasi antara 25\% dan 60\% sepanjang periode, dan Meksiko naik mendekati 100\% pada 1980. Grafik 3 mewakili "Pembangunan ekonomi"; ketiga negara bermula sekitar 25\% hingga 40\%. AS turun menjadi sekitar 5\% pada 1990 sebelum naik menjadi 20\%; Kanada turun menjadi sekitar 25\% pada 1985, naik menjadi 50\% pada 2000, lalu turun lagi menjadi 25\%; dan Meksiko naik menjadi sekitar 100\% pada 1980 sebelum turun menjadi sekitar 85\%. Grafik 4 mewakili "Hak asasi manusia"; semua negara mengelompok di sekitar 65\% pada 1945, lalu AS turun menjadi 25\% pada 1975 dan berfluktuasi antara 10\% dan 30\% selama sisa periode, Kanada perlahan turun menjadi sekitar 15\%, dan Meksiko naik mendekati 100\% pada 1985 lalu perlahan turun menjadi sekitar 80\%. Grafik 5 mewakili "Senjata dan bahan nuklir"; semua negara bermula mendekati 0\% pada 1945, lalu AS sedikit naik tetapi umumnya berfluktuasi antara 15\% dan 40\%, Kanada naik menjadi sekitar 60\% pada 1965 sebelum turun dan berfluktuasi di sekitar 40\% hingga 50\%, dan Meksiko naik cepat menjadi sekitar 90\% pada 1970 lalu mendekati 100\% seiring waktu. Grafik 6 mewakili "Konflik Palestina"; semua negara bermula antara 50\% dan 75\%, lalu AS terus turun menjadi sekitar 10\% pada 1985 dan setelah itu mendekati 5\%, Kanada sedikit turun menjadi sekitar 35\% pada 1970 sebelum naik menjadi sekitar 70\% pada 2000 dan kemudian turun cepat mendekati 0\%, sedangkan Meksiko naik bertahap menjadi sekitar 95\% pada 1985 lalu bertahan kurang lebih tetap.]{0.85}{eoce/unvotes}{unvotes}
}\par
\begin{parts}
\item Daftarkan variabel yang digunakan untuk membuat visualisasi ini.
\item Tentukan apakah setiap variabel dalam penelitian bersifat numerik atau
kategoris. Jika numerik, tentukan apakah variabel itu kontinu atau diskret.
Jika kategoris, nyatakan apakah variabel tersebut ordinal.
\end{parts}
}{}
}




%%%%%
\section{Prinsip dan strategi pengambilan sampel}
\label{overviewOfDataCollectionPrinciples}
\label{section_obs_data_sampling}

\index{sampel|(}
\index{populasi|(}

Langkah pertama dalam melakukan penelitian adalah menentukan topik
atau pertanyaan yang hendak diteliti.
Pertanyaan penelitian yang dirumuskan dengan jelas membantu menentukan
subjek atau kasus yang perlu dipelajari serta variabel yang penting.
Kita juga perlu mempertimbangkan \emph{bagaimana} data dikumpulkan
agar data tersebut andal dan membantu mencapai tujuan penelitian.


\subsection{Populasi dan sampel}
\label{populationsAndSamples}

\noindent%
Perhatikan tiga pertanyaan penelitian berikut:
\begin{enumerate}
\setlength{\itemsep}{0mm}
\item
    Berapakah rata-rata kadar merkuri pada ikan todak
    di Samudra Atlantik?
\item
    \label{timeToGraduationQuestionForUCLAStudents}%
    Selama 5 tahun terakhir, berapa rata-rata waktu yang dibutuhkan
    mahasiswa program sarjana Duke untuk menyelesaikan gelarnya?
\item
    \label{identifyPopulationOfStentStudy}%
    Apakah obat baru mengurangi jumlah kematian pada pasien
    dengan penyakit jantung berat?
\end{enumerate}
Setiap pertanyaan penelitian mengacu pada \term{populasi} sasaran. Pada
pertanyaan pertama, populasi sasarannya adalah seluruh ikan todak di Samudra
Atlantik, dan setiap ikan merupakan satu kasus. Mengumpulkan data untuk setiap
kasus dalam suatu populasi sering kali terlalu mahal. Sebagai gantinya, kita
mengambil sampel. Sebuah \term{sampel} merupakan sebagian dari kasus-kasus
tersebut dan biasanya hanya mencakup sebagian kecil populasi. Sebagai contoh,
60 ikan todak (atau jumlah lain) dapat dipilih dari populasi. Data sampel ini
kemudian dapat digunakan untuk memperkirakan rata-rata populasi dan menjawab
pertanyaan penelitian.

\begin{exercisewrap}
\begin{nexercise}\label{identifyingThePopulationForTwoQuestionsInPopAndSampSubsection}%
Untuk pertanyaan kedua dan ketiga di atas,
tentukan populasi sasaran dan apa yang
merupakan satu kasus individual.\footnotemark
\end{nexercise}
\end{exercisewrap}
\footnotetext{(\ref{timeToGraduationQuestionForUCLAStudents})
    Pertanyaan tentang waktu kelulusan hanya relevan bagi mahasiswa
    yang menyelesaikan gelarnya; rata-rata tersebut tidak dapat dihitung
    dengan memasukkan mahasiswa yang tidak pernah lulus. Karena itu,
    hanya mahasiswa program sarjana Duke yang lulus dalam lima tahun
    terakhir yang menjadi kasus dalam populasi yang sedang ditinjau.
    Setiap mahasiswa tersebut merupakan satu kasus individual.
    (\ref{identifyPopulationOfStentStudy})~Seseorang dengan penyakit
    jantung berat merupakan satu kasus. Populasinya mencakup semua orang
    dengan penyakit jantung berat.}


\subsection{Bukti anekdotal}
\label{anecdotalEvidenceSubsection}

\index{bias|(}

\noindent%
Perhatikan beberapa kemungkinan jawaban berikut untuk ketiga pertanyaan penelitian tadi:
\begin{enumerate}
\setlength{\itemsep}{0mm}
\item
    Seorang pria dalam berita mengalami keracunan merkuri setelah memakan
    ikan todak, sehingga rata-rata konsentrasi merkuri pada ikan todak pasti
    sangat tinggi dan berbahaya.
\item\label{iKnowThreeStudentsWhoTookMoreThan7YearsToGraduateAtDuke}
    Saya bertemu dua mahasiswa yang membutuhkan lebih dari 7 tahun untuk
    lulus dari Duke, sehingga menyelesaikan studi di Duke pasti memerlukan
    waktu lebih lama daripada di banyak perguruan tinggi lain.
\item\label{myFriendsDadDiedAfterSulphinpyrazon}
    Ayah teman saya mengalami serangan jantung dan meninggal setelah diberi
    obat baru untuk penyakit jantung, sehingga obat itu pasti tidak bekerja.
\end{enumerate}
Setiap kesimpulan tersebut didasarkan pada data.
Namun, terdapat dua masalah.
Pertama, data itu hanya mewakili satu atau dua kasus.
Kedua, dan yang lebih penting, belum jelas apakah kasus-kasus tersebut
benar-benar mewakili populasi.
Data yang dikumpulkan secara serampangan seperti ini disebut
\term{bukti anekdotal}.

\captionsetup{width=\textwidth-75mm}
\begin{figure}[h]
  \centering
  \hspace{8mm}\Figuress
    [Pemandangan musim dingin dengan pepohonan yang tertutup salju dan tumpukan salju besar di tepi jalan. Foto ini diambil di kampus University of Minnesota setelah sebuah badai; cabang-cabang pohon tampak putih terang karena tertutup salju.]
    {55mm}{mnWinter}{mnWinter}\hspace{4mm}
  \begin{minipage}[b]{\textwidth-75mm}
    \caption[bukti anekdotal]{Pada Februari 2010,
        sejumlah komentator media menggunakan badai salju besar di Amerika
        Serikat bagian timur laut untuk menentang pemanasan global. Jon Stewart
        menyindir penalaran tersebut: satu kejadian cuaca setempat tidak dapat
        mewakili pola iklim global jangka panjang. Foto: David M. Diez,
        CC BY-SA 3.0.
    \label{mnWinter}}
  \end{minipage}
\end{figure}
\captionsetup{width=\mycaptionwidth}

\begin{onebox}{Bukti anekdotal}
Waspadai data yang dikumpulkan secara serampangan.
Bukti semacam itu mungkin benar dan dapat diverifikasi, tetapi mungkin
hanya mewakili kasus-kasus luar biasa.
\end{onebox}

\D{\newpage}

Bukti anekdotal biasanya terdiri atas kasus-kasus tidak lazim yang kita ingat
karena cirinya menonjol. Sebagai contoh, kita lebih mungkin mengingat dua orang
yang kita temui dan membutuhkan 7~tahun untuk lulus daripada enam orang lain
yang lulus dalam empat tahun. Alih-alih hanya melihat kasus yang paling tidak
lazim, kita sebaiknya memeriksa sampel yang terdiri atas banyak kasus dan
mewakili populasi.

\subsection{Mengambil sampel dari populasi}

\index{sampel!sampel acak|(}
\index{sampel!bias|(}

Kita dapat mencoba memperkirakan waktu kelulusan mahasiswa program sarjana
Duke selama 5 tahun terakhir dengan mengumpulkan sampel mahasiswa.
Semua lulusan dalam 5 tahun terakhir merupakan
\emph{populasi}\index{populasi}, sedangkan lulusan yang dipilih untuk
ditinjau secara bersama-sama disebut
\emph{sampel}\index{sampel}.
Secara umum, kita selalu berupaya memilih sampel dari suatu populasi
\emph{secara acak}.
Bentuk pemilihan acak yang paling dasar setara dengan cara melakukan undian.
Sebagai contoh, untuk memilih lulusan, kita dapat menuliskan nama setiap
lulusan pada kupon undian lalu mengambil 100 kupon.
Nama-nama yang terambil akan membentuk sampel acak berisi 100 lulusan.
Kita memilih sampel secara acak untuk mengurangi kemungkinan timbulnya bias.

\begin{figure}[ht]
  \centering
  \Figures
    [Diagram memperlihatkan lingkaran besar di kiri untuk "semua lulusan" dan lingkaran kecil di kanan untuk "sampel". Banyak titik tersebar secara acak di dalam lingkaran kiri. Lima titik di antaranya memiliki panah menuju lima titik di dalam lingkaran kanan; tidak ada titik lain dalam lingkaran kanan.]
    {0.5}{popToSample}{popToSampleGraduates}
  \caption{Pada diagram ini, lima lulusan dipilih secara acak
      dari populasi untuk dimasukkan ke dalam sampel.}
  \label{popToSampleGraduates}
\end{figure}

\begin{examplewrap}
\begin{nexample}{Misalkan kita meminta seorang mahasiswa jurusan gizi
    memilih beberapa lulusan untuk penelitian ini.
    Menurut Anda, mahasiswa seperti apa yang mungkin dipilihnya?
    Apakah sampelnya akan mewakili semua lulusan?}
  Ia mungkin memilih lulusan dari bidang kesehatan dalam proporsi yang
  terlalu besar. Atau, pilihannya mungkin saja mewakili populasi dengan baik.
  Ketika memilih sampel secara manual, kita berisiko memperoleh sampel yang
  \termsub{bias}{sampel!bias}, sekalipun bias itu tidak disengaja.
\end{nexample}
\end{examplewrap}

\begin{figure}
  \centering
  \Figures
    [Diagram memperlihatkan lingkaran besar di kiri untuk "semua lulusan" dan lingkaran kecil di kanan untuk "sampel". Banyak titik tersebar secara acak di lingkaran kiri. Di dalamnya terdapat lingkaran lebih kecil berlabel "lulusan bidang kesehatan" yang memuat sebagian titik. Lima titik dari lingkaran kecil tersebut memiliki panah menuju lima titik dalam lingkaran sampel; tidak ada titik lain dalam lingkaran sampel.]
    {0.5}{popToSample}{popToSubSampleGraduates}
  \caption{Ketika diminta memilih sampel lulusan,
      seorang mahasiswa jurusan gizi tanpa sengaja mungkin memilih
      lulusan dari bidang kesehatan dalam proporsi yang terlalu besar.}
  \label{popToSubSampleGraduates}
\end{figure}

Jika seseorang bebas menentukan sendiri lulusan mana yang masuk ke dalam
sampel, sampel itu sangat mungkin condong ke minat orang tersebut, meskipun
hal itu sama sekali tidak disengaja.
Keadaan ini memasukkan \term{bias} ke dalam sampel.
Pengambilan sampel secara acak membantu mengatasi masalah ini.
Sampel acak yang paling dasar disebut
\term{sampel acak sederhana}; prosedurnya setara dengan menggunakan
undian untuk memilih kasus. Untuk ukuran sampel yang telah ditetapkan,
setiap himpunan kasus yang mungkin dengan ukuran tersebut memiliki peluang
yang sama untuk terpilih.

Mengambil sampel acak sederhana membantu meminimalkan bias.
Namun, bias dapat muncul melalui cara lain.
Bahkan ketika orang dipilih secara acak, misalnya untuk suatu survei,
kita perlu berhati-hati jika
\term{tingkat nonrespons}
\index{sampel!tingkat nonrespons|textbf} tinggi.
Sebagai contoh, jika hanya 30\% orang yang dipilih secara acak untuk suatu
survei benar-benar merespons, belum jelas apakah hasilnya
\term{representatif} bagi seluruh populasi.
\term{Bias nonrespons}
\index{sampel!bias nonrespons|textbf} ini dapat menyimpangkan hasil.

\begin{figure}[h]
  \centering
  \Figures
    [Diagram memperlihatkan lingkaran besar di kiri untuk "populasi sasaran" dan lingkaran kecil di kanan untuk "sampel". Banyak titik tersebar secara acak di dalam lingkaran kiri. Di dalamnya terdapat lingkaran lebih kecil berlabel "populasi yang benar-benar tersampel" yang memuat sebagian titik. Lima titik dari lingkaran kecil tersebut memiliki panah menuju lima titik dalam lingkaran sampel; tidak ada titik lain dalam lingkaran sampel.]
    {0.5}{popToSample}{surveySample}
  \caption{Karena sebagian orang mungkin tidak merespons,
      suatu survei mungkin hanya menjangkau kelompok tertentu
      dalam populasi. Masalah ini sulit, dan sering kali mustahil,
      diperbaiki sepenuhnya.}
  \label{surveySample}
\end{figure}

Kelemahan umum lainnya adalah
\term{sampel kemudahan}\index{sampel!sampel kemudahan},
yakni individu yang mudah dijangkau lebih mungkin dimasukkan ke dalam sampel.
Sebagai contoh, jika survei politik dilakukan dengan menghentikan pejalan kaki
di Bronx, sampelnya tidak akan mewakili seluruh Kota New York.
Sering kali sulit menentukan subpopulasi mana yang sebenarnya diwakili oleh
sampel kemudahan.

\begin{exercisewrap}
\begin{nexercise}
Kita dapat dengan mudah mengakses penilaian terhadap produk, penjual, dan
perusahaan melalui situs web. Penilaian ini hanya berasal dari orang-orang
yang secara khusus meluangkan waktu untuk memberikannya. Jika 50\% ulasan
daring terhadap suatu produk bersifat negatif, apakah menurut Anda itu berarti
50\% pembeli tidak puas dengan produk tersebut?\footnotemark
\end{nexercise}
\end{exercisewrap}
\footnotetext{Jawaban dapat beragam.
    Berdasarkan pengalaman anekdotal kami sendiri, orang tampaknya lebih
    cenderung mengeluhkan produk yang tidak memenuhi harapan daripada memuji
    produk yang bekerja sesuai harapan. Karena itu, kami menduga terdapat bias
    negatif dalam penilaian produk di situs seperti Amazon. Namun, karena
    pengalaman kami mungkin tidak representatif, kami tetap terbuka terhadap
    kemungkinan lain.}

\index{sampel!bias|)}
\index{sampel!sampel acak|)}
\index{bias|)}
\index{populasi|)}
\index{sampel|)}

\subsection{Studi observasional}

Data yang dikumpulkan tanpa secara eksplisit memberikan
(atau menahan) perlakuan disebut \term{data observasional}.
Sebagai contoh, data pinjaman dan data county yang dijelaskan pada
Bagian~\ref{dataBasics}
keduanya merupakan contoh data observasional.
%It is important to collect such data in
%a thoughtful and rigorous manner so that statistical
%analyses based on the data can have meaningful
%and generalizable results.
Menarik kesimpulan kausal dari eksperimen yang dirancang dengan baik sering
kali masuk akal. Namun, menarik kesimpulan kausal yang sama dari data
observasional dapat menyesatkan dan tidak dianjurkan.
Karena itu, studi observasional umumnya hanya memadai untuk menunjukkan
asosiasi atau membentuk hipotesis yang kemudian kita periksa melalui
eksperimen.

\begin{exercisewrap}
\begin{nexercise}\label{sunscreenLurkingExample}%
Misalkan sebuah studi observasional mengikuti penggunaan tabir surya dan
kanker kulit, lalu menemukan bahwa makin banyak tabir surya yang digunakan
seseorang, makin besar kemungkinan orang tersebut menderita kanker kulit.
Apakah ini berarti tabir surya \emph{menyebabkan} kanker kulit?\footnotemark
\end{nexercise}
\end{exercisewrap}
\footnotetext{Tidak.
    Lihat paragraf setelah latihan untuk penjelasannya.}

Sejumlah penelitian terdahulu menunjukkan bahwa penggunaan tabir surya justru
menurunkan risiko kanker kulit. Karena itu, mungkin terdapat variabel lain yang
dapat menjelaskan asosiasi hipotetis antara penggunaan tabir surya dan kanker
kulit ini. Salah satu informasi penting yang tidak disertakan adalah paparan
sinar matahari. Seseorang yang berada di bawah sinar matahari sepanjang hari
lebih mungkin menggunakan tabir surya \emph{dan} lebih mungkin terkena kanker
kulit. Penyelidikan sederhana tadi tidak memperhitungkan paparan sinar matahari.
\begin{center}
  \Figures
    [Tiga kotak berisi teks tersusun membentuk segitiga. Kotak "paparan sinar matahari" memiliki dua panah yang masing-masing menuju kotak "menggunakan tabir surya" dan "kanker kulit". Panah ketiga berwarna lebih muda mengarah dari kotak "menggunakan tabir surya" ke kotak "kanker kulit", dengan tanda tanya di atas panah tersebut.]
    {0.55}{variables}{sunCausesCancer}
\end{center}
% Some studies:
% http://www.sciencedirect.com/science/article/pii/S0140673698121682
% http://archderm.ama-assn.org/cgi/content/abstract/122/5/537
% Study with a similar scenario to that described here:
% http://onlinelibrary.wiley.com/doi/10.1002/ijc.22745/full

Paparan sinar matahari disebut \term{variabel perancu},\footnote{Disebut juga
\term{variabel tersembunyi}, \term{faktor perancu}, atau \term{perancu}.}
yaitu variabel yang berkorelasi dengan variabel penjelas sekaligus variabel
respons. Salah satu cara untuk mencoba membenarkan kesimpulan kausal dari studi
observasional adalah mencari variabel perancu selengkap mungkin. Namun, tidak
ada jaminan bahwa semua variabel perancu dapat diperiksa atau diukur.

%In the same way, the \data{county} data set is an observational study with confounding variables, and its data cannot easily be used to make causal conclusions.

\begin{exercisewrap}
\begin{nexercise}
Gambar~\ref{multiunitsVsOwnership} menunjukkan asosiasi negatif
antara tingkat kepemilikan rumah dan persentase bangunan multiunit
di suatu county. Namun, menyimpulkan bahwa kedua variabel tersebut memiliki
hubungan kausal tidaklah beralasan.
Usulkan sebuah variabel yang mungkin menjelaskan hubungan
negatif tersebut.\footnotemark
\end{nexercise}
\end{exercisewrap}
\footnotetext{Jawaban dapat beragam.
    Kepadatan penduduk mungkin berperan penting. Jika suatu county sangat
    padat, sebagian lebih besar penduduknya mungkin perlu tinggal di bangunan
    multiunit. Selain itu, kepadatan tinggi dapat turut menaikkan nilai properti,
    sehingga kepemilikan rumah tidak terjangkau bagi banyak penduduk.}

Dua bentuk studi observasional yang umum dan dibahas di sini adalah
studi prospektif dan studi retrospektif.
Sebuah \term{studi prospektif} mengidentifikasi individu
dan mengumpulkan informasi seiring berlangsungnya peristiwa.
Sebagai contoh, peneliti medis dapat mengidentifikasi dan mengikuti
sekelompok pasien selama bertahun-tahun untuk menilai kemungkinan pengaruh
perilaku terhadap risiko kanker.
Salah satu contohnya adalah \emph{The Nurses' Health Study},
yang dimulai pada 1976 dan diperluas pada 1989.
Studi prospektif ini merekrut perawat terdaftar, kemudian mengumpulkan data
dari mereka melalui kuesioner.
\termsub{Studi retrospektif}{studi retrospektif}
mengumpulkan data setelah peristiwa terjadi; misalnya, peneliti dapat
menelaah peristiwa masa lalu dalam rekam medis.
Sejumlah himpunan data dapat memuat variabel yang dikumpulkan secara
prospektif maupun retrospektif.


\subsection{Empat metode pengambilan sampel}
\label{fourSamplingMethods}
\label{threeSamplingMethods}

Hampir semua metode statistik didasarkan pada gagasan bahwa terdapat keacakan
yang mendasari data. Jika data observasional dari suatu populasi tidak
dikumpulkan dalam kerangka acak, metode-metode statistik tersebut---baik
estimasi maupun galat yang menyertai estimasi---tidak dapat diandalkan. Di sini
kita membahas empat teknik pengambilan sampel acak: pengambilan sampel acak
sederhana, berstrata, klaster, dan multitahap. Gambar~\ref{simple_stratified}
dan~\ref{cluster_multistage} menyajikan gambaran grafis teknik-teknik tersebut.

\begin{figure}
  \centering
  \Figures
    [Dua diagram ditampilkan bertumpuk. Diagram atas berupa persegi panjang besar yang memuat banyak titik; 18 titik dilingkari dan diberi warna berbeda. Diagram bawah juga berupa persegi panjang besar, tetapi memuat 6 lingkaran besar berlabel "Strata 1" hingga "Strata 6". Setiap lingkaran memuat banyak titik, dan pada masing-masing strata terdapat 3 titik yang dilingkari serta diberi warna berbeda, serupa dengan penandaan 18 titik pada diagram atas.]
    {}{samplingMethodsFigure}{simple_stratified}
  \caption{
      Contoh pengambilan sampel acak sederhana
      \index{sampel!pengambilan sampel acak sederhana}
      dan pengambilan sampel berstrata.
      \index{sampel!pengambilan sampel berstrata}
      Pada panel atas, pengambilan sampel acak sederhana digunakan untuk
      memilih 18 kasus secara acak. Pada panel bawah, digunakan pengambilan
      sampel berstrata: kasus dikelompokkan ke dalam strata, lalu pengambilan
      sampel acak sederhana diterapkan di \mbox{setiap strata}.}
  \label{simple_stratified}
\end{figure}

\termsub{Pengambilan sampel acak sederhana}{sampel!pengambilan sampel acak sederhana}
mungkin merupakan bentuk pengambilan sampel acak yang paling intuitif.
Perhatikan gaji pemain Major League Baseball (MLB), dengan setiap pemain
tergabung dalam salah satu dari 30 tim liga. Untuk mengambil sampel acak
sederhana yang terdiri atas 120 pemain beserta gajinya, kita dapat menuliskan
nama ratusan pemain pada musim tersebut di secarik kertas, memasukkan semua
kertas ke dalam wadah, mengocoknya hingga seluruh nama tercampur, lalu
mengambil kertas sampai terkumpul sampel 120 pemain. Secara umum, sebuah
sampel disebut ``acak sederhana'' jika setiap kasus dalam populasi memiliki
peluang yang sama untuk masuk ke dalam sampel akhir \emph{dan} informasi bahwa
satu kasus masuk ke dalam sampel tidak memberikan informasi berguna mengenai
kasus lain mana yang turut masuk.

\termsub{Pengambilan sampel berstrata}{sampel!pengambilan sampel berstrata}
merupakan strategi pengambilan sampel yang membagi masalah menjadi bagian-bagian.
Populasi dibagi ke dalam kelompok yang disebut
\term{strata}\index{sampel!strata|textbf}.
Strata dipilih agar kasus-kasus yang serupa dikelompokkan bersama, lalu metode
pengambilan sampel kedua---biasanya pengambilan sampel acak sederhana---
diterapkan dalam setiap strata.
Pada contoh gaji pemain bisbol, tim dapat menjadi strata karena sebagian tim
memiliki dana jauh lebih besar (hingga 4~kali lipat!).
Kita kemudian dapat memilih 4 pemain secara acak dari setiap tim sehingga
jumlah keseluruhannya 120 pemain.

Pengambilan sampel berstrata sangat berguna ketika kasus-kasus dalam setiap
strata sangat serupa dalam hal luaran yang menjadi perhatian. Kekurangannya,
menganalisis data dari sampel berstrata lebih rumit daripada menganalisis data
dari sampel acak sederhana. Metode analisis yang diperkenalkan dalam buku ini
perlu diperluas untuk menganalisis data yang dikumpulkan melalui pengambilan
sampel berstrata.

\begin{examplewrap}
\begin{nexample}{Mengapa keserupaan yang tinggi antarkasus
    dalam setiap strata menguntungkan?}
  Jika kasus-kasus sangat serupa, kita dapat memperoleh estimasi yang lebih
  stabil untuk subpopulasi dalam suatu strata, sehingga estimasi pada setiap
  kelompok menjadi lebih presisi. Ketika estimasi-estimasi ini digabungkan
  menjadi satu estimasi bagi seluruh populasi, estimasi populasi tersebut
  cenderung lebih presisi karena setiap estimasi kelompok juga lebih presisi.
\end{nexample}
\end{examplewrap}

Dalam \termsub{sampel klaster}{sampel!sampel klaster}, populasi dibagi menjadi
banyak kelompok yang disebut \termsub{klaster}{sampel!klaster}. Sejumlah
klaster yang telah ditetapkan kemudian dipilih, dan semua observasi dalam
setiap klaster terpilih dimasukkan ke dalam sampel. Sebuah
\termsub{sampel multitahap}{sampel!sampel multitahap} menyerupai sampel klaster,
tetapi alih-alih menyertakan semua observasi dalam setiap klaster, kita
mengambil sampel acak di dalam setiap klaster terpilih. %Multistage sampling is similar to stratified sampling in its process, except that stratified sampling requires observations be sampled from \emph{every} stratum.

\begin{figure}
  \centering
  \Figures
    [Dua diagram ditampilkan bertumpuk. Diagram atas berupa persegi panjang besar yang memuat 9 lingkaran besar berlabel "Klaster 1" hingga "Klaster 9". Semua lingkaran besar memuat titik. Tiga lingkaran---Klaster 3, Klaster 4, dan Klaster 8---serta seluruh titik di dalamnya diberi warna berbeda. Diagram bawah serupa, tetapi hanya 6 titik dalam masing-masing dari 3 klaster terpilih yang diberi warna berbeda.]
    {}{samplingMethodsFigure}{cluster_multistage}
  \caption{Contoh pengambilan sampel klaster
  \index{sampel!pengambilan sampel klaster}
  dan multitahap\index{sampel!pengambilan sampel multitahap}.
  Pada panel atas, digunakan pengambilan sampel klaster:
  data dibagi ke dalam sembilan klaster, tiga klaster dipilih sebagai sampel,
  dan semua observasi dalam ketiga klaster tersebut dimasukkan ke dalam sampel.
  Pada panel bawah, digunakan pengambilan sampel multitahap. Perbedaannya dari
  pengambilan sampel klaster adalah bahwa kita memilih secara acak sebagian
  anggota setiap klaster untuk dimasukkan ke dalam sampel, alih-alih mengukur
  setiap kasus dalam tiap klaster terpilih.}
\label{cluster_multistage}
\end{figure}

Pengambilan sampel klaster atau multitahap kadang-kadang lebih ekonomis daripada
teknik pengambilan sampel lainnya.
Selain itu, tidak seperti pengambilan sampel berstrata, kedua pendekatan ini
paling berguna ketika terdapat banyak keragaman antarkasus di dalam setiap
klaster, sedangkan klaster-klaster itu sendiri tidak terlalu berbeda satu sama
lain. Sebagai contoh, jika lingkungan tempat tinggal menjadi klaster,
pengambilan sampel klaster atau multitahap bekerja paling baik ketika setiap
lingkungan sangat beragam.
Kekurangan metode-metode ini adalah bahwa analisis datanya biasanya memerlukan
teknik yang lebih lanjut, meskipun metode dalam buku ini dapat diperluas untuk
menangani data semacam itu.

\begin{examplewrap}
\begin{nexample}{Misalkan kita hendak memperkirakan tingkat malaria
    di suatu wilayah pedesaan berhutan tropis lebat di Indonesia.
    Kita mengetahui bahwa terdapat 30 desa di bagian kawasan hutan tersebut,
    dan setiap desa kurang lebih serupa dengan desa lainnya.
    Tujuan kita adalah melakukan tes malaria pada 150 orang.
    Metode pengambilan sampel apa yang sebaiknya digunakan?}
  Sampel acak sederhana kemungkinan akan memilih orang dari seluruh 30 desa,
  sehingga pengumpulan data bisa menjadi sangat mahal.
  Pengambilan sampel berstrata akan sulit karena tidak jelas bagaimana kita
  dapat membentuk strata berisi individu yang serupa.
  Namun, pengambilan sampel klaster atau multitahap tampak sangat sesuai.
  Jika memutuskan menggunakan pengambilan sampel multitahap, kita dapat memilih
  separuh desa secara acak, lalu memilih 10 orang secara acak dari setiap desa.
  Cara ini kemungkinan besar akan mengurangi biaya pengumpulan data secara
  berarti dibandingkan sampel acak sederhana, dan sampel multitahap tersebut
  tetap akan memberi kita informasi yang andal, meskipun analisis datanya
  memerlukan metode yang sedikit lebih lanjut daripada yang dibahas dalam
  buku ini.
\end{nexample}
\end{examplewrap}


{\begingroup\let\clearpageforsection\relax\exercisesheader{}\endgroup

% 13

\eoce{\qt{Pencemaran udara dan luaran kelahiran: cakupan inferensi\label{scope_airpoll}} 
Latihan~\ref{study_components_airpoll} memperkenalkan sebuah studi yang 
mengumpulkan data untuk menelaah hubungan antara pencemar udara dan kelahiran 
prematur di California Selatan. Selama studi, tingkat pencemaran udara diukur 
oleh stasiun pemantau kualitas udara. Data lama kehamilan dikumpulkan untuk 
143.196 kelahiran antara 1989 dan 1993, dan paparan pencemaran udara selama 
kehamilan dihitung untuk setiap kelahiran.
\begin{parts}
\item Tentukan populasi sasaran dan sampel dalam studi ini.
\item Jelaskan apakah hasil studi dapat digeneralisasikan ke populasi dan apakah 
temuannya dapat digunakan untuk menetapkan hubungan kausal.
\end{parts}
}{}

% 14

\eoce{\qt{Kecurangan: cakupan inferensi\label{scope_cheaters}} 
Latihan~\ref{study_components_cheaters} memperkenalkan studi tentang hubungan 
antara kejujuran, usia, dan pengendalian diri. Para peneliti melakukan eksperimen 
pada 160 anak berusia 5 hingga 15 tahun. Setiap anak diminta melempar koin 
seimbang tanpa dilihat orang lain dan mencatat hasilnya (putih atau hitam) pada 
selembar kertas; hanya anak yang melaporkan hasil putih yang akan diberi hadiah. 
Separuh anak 
secara eksplisit diminta untuk tidak berbuat curang, sedangkan separuh lainnya 
tidak diberi instruksi eksplisit. Terdapat perbedaan tingkat kecurangan antara 
kelompok dengan dan tanpa instruksi, serta sejumlah perbedaan menurut 
karakteristik anak di dalam setiap kelompok.
\begin{parts}
\item Tentukan populasi sasaran dan sampel dalam studi ini.
\item Jelaskan apakah hasil studi dapat digeneralisasikan ke populasi dan apakah 
temuannya dapat digunakan untuk menetapkan hubungan kausal.
\end{parts}
}{}

% 15

\eoce{\qt{Metode Buteyko: cakupan inferensi\label{scope_buteyko}} 
Latihan~\ref{study_components_buteyko} memperkenalkan studi mengenai penggunaan 
teknik pernapasan dangkal Buteyko untuk mengurangi gejala asma dan meningkatkan 
kualitas hidup. Studi ini merekrut 600 pasien asma berusia 18--69 tahun yang 
mengandalkan obat untuk perawatan asma, lalu mengalokasikan mereka secara acak 
ke dua kelompok: satu kelompok mempraktikkan metode Buteyko dan kelompok lainnya 
tidak. Kelompok Buteyko rata-rata mengalami penurunan gejala asma yang signifikan 
dan peningkatan kualitas hidup.
\begin{parts}
\item Tentukan populasi sasaran dan sampel dalam studi ini.
\item Jelaskan apakah hasil studi dapat digeneralisasikan ke populasi dan apakah 
temuannya dapat digunakan untuk menetapkan hubungan kausal.
\end{parts}
}{}

% 16

\eoce{\qt{Mengambil permen: cakupan inferensi\label{scope_stealers}} 
Latihan~\ref{study_components_stealers} memperkenalkan studi tentang hubungan 
antara kelas sosial ekonomi dan perilaku tidak etis. Dalam studi ini, 129 
mahasiswa program sarjana University of California, Berkeley diminta 
menggolongkan diri ke kelas sosial rendah atau tinggi dengan membandingkan diri 
mereka dengan orang yang memiliki paling banyak (atau paling sedikit) uang dan 
pendidikan, serta pekerjaan yang paling (atau paling tidak) dihormati. Mereka 
juga diberi sebuah stoples berisi permen yang dibungkus satu per satu. Peneliti 
memberi tahu bahwa permen tersebut ditujukan bagi anak-anak di laboratorium 
terdekat, tetapi peserta boleh mengambil beberapa jika ingin. Setelah 
menyelesaikan tugas lain yang tidak berkaitan, peserta melaporkan jumlah permen 
yang mereka ambil. Peserta yang tergolong kelas atas ditemukan mengambil lebih 
banyak permen daripada peserta lainnya.
\begin{parts}
\item Tentukan populasi sasaran dan sampel dalam studi ini.
\item Jelaskan apakah hasil studi dapat digeneralisasikan ke populasi dan apakah 
temuannya dapat digunakan untuk menetapkan hubungan kausal.
\end{parts}
}{}

% 17

\eoce{\qt{Bersantai setelah bekerja\label{relax_after_work_definitions}} General 
Social Survey mengambil sampel acak 1.155 warga Amerika dan menanyakan berapa 
jam yang mereka miliki setelah hari kerja biasa untuk bersantai atau melakukan 
kegiatan yang mereka sukai. Rata-rata waktu bersantai yang diperoleh adalah 
1,65 jam. Tentukan mana di antara hal-hal berikut yang merupakan observasi, 
variabel, statistik sampel (nilai yang dihitung berdasarkan sampel teramati), 
atau parameter populasi.
\begin{parts}
\item Seorang warga Amerika dalam sampel.
\item Jumlah jam yang digunakan untuk bersantai setelah hari kerja biasa.
\item 1,65.
\item Rata-rata jumlah jam yang digunakan seluruh warga Amerika untuk bersantai 
setelah hari kerja biasa.
\end{parts}
}{}

% 18

\eoce{\qt{Kucing di YouTube\label{cats_on_youtube_definitions}} Misalkan Anda 
ingin memperkirakan persentase video YouTube yang menampilkan kucing. Karena 
mustahil menonton semua video di YouTube, Anda menggunakan pemilih video acak 
untuk memilih 1.000 video. Anda menemukan bahwa 2\% dari video tersebut 
menampilkan kucing. Tentukan mana di antara hal-hal berikut yang merupakan 
observasi, variabel, statistik sampel (nilai yang dihitung berdasarkan sampel 
teramati), atau parameter populasi.
\begin{parts}
\item Persentase seluruh video YouTube yang menampilkan kucing.
\item 2\%.
\item Sebuah video dalam sampel Anda.
\item Apakah sebuah video menampilkan kucing atau tidak.
\end{parts}
}{}

% 19

\eoce{\qt{Kepuasan mata kuliah antarkelompok\label{course_satisfaction_sections}} 
Sebuah kelas besar di perguruan tinggi memiliki 160 mahasiswa. Seluruh 
mahasiswa mengikuti kuliah bersama, tetapi untuk sesi praktikum mereka dibagi 
menjadi 4 kelompok yang masing-masing berisi 40 mahasiswa dan dibimbing oleh 
asisten pengajar yang berbeda. Dosen hendak melakukan survei mengenai kepuasan 
mahasiswa terhadap mata kuliah tersebut. Ia menduga kelompok praktikum yang 
diikuti mahasiswa dapat memengaruhi kepuasan mereka secara keseluruhan.
\begin{parts}
\item Jenis studi apakah ini?
\item Usulkan strategi pengambilan sampel untuk melaksanakan studi ini.
\end{parts}
}{}

% 20

\eoce{\qt{Usulan hunian di berbagai asrama\label{housing_proposal_dorms}} Di sebuah 
kampus besar, mahasiswa tahun pertama dan kedua tinggal di asrama di bagian 
timur kampus, sedangkan mahasiswa tahun ketiga dan keempat tinggal di asrama 
di bagian barat. Misalkan Anda ingin mengumpulkan pendapat mahasiswa tentang 
struktur hunian baru yang diusulkan pengelola perguruan tinggi dan ingin 
memastikan bahwa survei mewakili pendapat mahasiswa dari setiap angkatan 
secara setara.
\begin{parts}
\item Jenis studi apakah ini?
\item Usulkan strategi pengambilan sampel untuk melaksanakan studi ini.
\end{parts}
}{}

% 21

\eoce{\qt{Penggunaan internet dan harapan hidup\label{internet_life_expectancy}} 
Diagram pencar berikut dibuat dalam sebuah studi yang menilai hubungan antara 
estimasi harapan hidup saat lahir (per 2014) dan persentase pengguna internet 
(per 2009) di 208 negara yang memiliki kedua jenis data tersebut.\footfullcite{data:ciaFactbook}

\noindent\begin{minipage}[c]{0.44\textwidth}
\begin{parts}
\item Jelaskan hubungan antara harapan hidup dan persentase pengguna internet.
\item Jenis studi apakah ini?
\item Sebutkan kemungkinan variabel perancu yang dapat menjelaskan hubungan ini 
dan uraikan kemungkinan pengaruhnya.
\end{parts} \vspace{15mm}
\end{minipage}
\begin{minipage}[r]{0.55\textwidth}
\hfill%
\Figures[Diagram pencar dengan "persentase pengguna internet" (0\% hingga 100\%) pada sumbu horizontal dan "harapan hidup saat lahir" (50 hingga 90) pada sumbu vertikal. Untuk 0\% hingga 15\%, sekitar 100 titik tersebar cukup merata antara 50 dan 75. Untuk 15\% hingga 90\%, titik-titik terkonsentrasi antara sekitar 70 dan 85 serta menunjukkan sedikit tren menaik.]{0.87}{eoce/internet_life_expectancy}{internet_life_expectancy}
\end{minipage}
}{}

% 22

\eoce{\qt{Dilanda stres, Bagian I\label{stressed_out_observational}} Sebuah studi 
yang menyurvei sampel acak siswa sekolah menengah yang secara umum sehat 
menemukan bahwa mereka lebih mungkin mengalami kram otot ketika stres. Studi 
tersebut juga mencatat bahwa siswa minum lebih banyak kopi dan tidur lebih 
sedikit ketika stres.
\begin{parts}
\item Jenis studi apakah ini?
\item Dapatkah studi ini digunakan untuk menyimpulkan hubungan kausal antara 
peningkatan stres dan kram otot?
\item Sebutkan kemungkinan variabel perancu yang dapat menjelaskan hubungan 
teramati antara peningkatan stres dan kram otot. 
\end{parts}
}{}

% 23

\eoce{\qt{Evaluasi metode pengambilan sampel\label{evaluate_sampling_methods}} 
Sebuah universitas ingin mengetahui proporsi mahasiswa program sarjananya yang 
mendukung iuran tahunan baru sebesar \$25 untuk memperbaiki pusat kegiatan 
mahasiswa. Untuk setiap metode yang diusulkan berikut, nyatakan apakah metode 
tersebut masuk akal atau tidak.
\begin{parts}
\item Survei sampel acak sederhana yang terdiri atas 500 mahasiswa.
\item Bagi mahasiswa ke dalam strata berdasarkan bidang studi, lalu ambil sampel 
10\% mahasiswa dari setiap strata.
\item Kelompokkan mahasiswa ke dalam klaster berdasarkan usia (misalnya usia 
18 tahun dalam satu klaster, usia 19 tahun dalam satu klaster, dan seterusnya), 
lalu pilih tiga klaster secara acak dan survei seluruh mahasiswa dalam klaster 
tersebut.
\end{parts}
}{}

% 24

\eoce{\qt{Pemanggilan digit acak\label{random_digit_dialing}} Gallup Poll 
menggunakan prosedur yang disebut pemanggilan digit acak. Prosedur ini membuat 
nomor telepon berdasarkan daftar seluruh kode area di Amerika Serikat serta 
jumlah rumah tangga yang terkait dengan setiap kode area. Berikan satu alasan 
yang mungkin menjelaskan mengapa Gallup Poll menggunakan pemanggilan digit acak 
alih-alih memilih nomor telepon dari buku telepon.
}{}

% 25

\eoce{\qt{Kecenderungan menyukai atau tidak menyukai\label{scope_haters}}
Sebuah studi yang diterbitkan dalam
\textit{Journal of Personality and Social Psychology}
meminta 200 peserta dari kalangan laki-laki dan perempuan yang dipilih secara
acak untuk menilai
perasaan mereka terhadap berbagai topik, seperti berkemah, layanan kesehatan,
arsitektur, taksidermi, teka-teki silang, dan Jepang. Tujuannya adalah mengukur
sikap mereka terhadap rangsangan yang sebagian besar tidak saling berkaitan.
Selanjutnya, peserta menerima informasi tentang produk baru berupa oven
gelombang mikro. Produk ini sebenarnya tidak ada, tetapi peserta tidak
mengetahuinya; mereka diberi tiga ulasan positif dan tiga ulasan negatif yang
direkayasa. Orang yang bereaksi positif terhadap topik-topik pada pengukuran
sikap disposisional juga cenderung bereaksi positif terhadap oven tersebut,
sedangkan orang yang bereaksi negatif cenderung menilainya secara negatif.
Para peneliti menafsirkan pola itu sebagai kecenderungan evaluatif umum:
sebagian peserta cenderung menilai banyak objek secara lebih positif, sedangkan
peserta lain cenderung menilainya secara lebih negatif. Menurut mereka, pola
lintas-objek ini dapat membantu menjelaskan bagaimana sikap terbentuk.
\footfullcite{Hepler:2013}
\begin{parts}
\item Apa saja yang menjadi kasus?
\item Apa variabel respons dalam studi ini?
\item Apa variabel penjelas dalam studi ini?
\item Apakah studi ini menggunakan pengambilan sampel acak?
\item Apakah ini studi observasional atau eksperimen? Jelaskan alasan Anda.
\item Dapatkah kita menetapkan hubungan kausal antara variabel penjelas dan 
variabel respons?
\item Dapatkah hasil studi digeneralisasikan ke populasi secara luas?
\end{parts}
}{}

% 26

\eoce{\qt{Ukuran keluarga\label{family_size}} Misalkan kita ingin memperkirakan 
ukuran rumah tangga. Dalam konteks ini, rumah tangga terdiri atas orang-orang 
yang tinggal bersama dalam hunian yang sama dan berbagi fasilitas tempat 
tinggal. Jika kita memilih murid secara acak di sebuah sekolah dasar dan 
menanyakan ukuran keluarga mereka, apakah hasilnya akan mengukur ukuran rumah 
tangga dengan baik? Atau apakah rata-ratanya akan bias? Jika bias, apakah 
hasilnya akan mengestimasi nilai sebenarnya terlalu tinggi atau terlalu rendah?
}{}

% 27

\eoce{\qt{Strategi pengambilan sampel\label{sampling_strategies}} Seorang mahasiswa 
statistika ingin mengetahui hubungan antara waktu yang dihabiskan siswa di situs 
jejaring sosial dan kinerja mereka di sekolah, lalu memutuskan untuk melakukan 
survei. Beberapa strategi penelitian untuk mengumpulkan data diuraikan berikut. 
Untuk setiap strategi, sebutkan metode pengambilan sampel yang diusulkan dan 
bias yang mungkin timbul.
\begin{parts}
\item Ia memilih 40 siswa secara acak dari populasi studi, memberikan survei 
kepada mereka, lalu meminta mereka mengisinya dan mengembalikannya keesokan hari.
\item Ia hanya membagikan survei kepada teman-temannya dan memastikan setiap 
teman mengisinya.
\item Ia mengunggah tautan survei daring di Facebook dan meminta teman-temannya 
mengisi survei tersebut.
\item Ia memilih 5 kelas secara acak, lalu meminta sampel acak siswa dari 
kelas-kelas tersebut untuk mengisi survei.
\end{parts}
}{}

% 28

\eoce{\qt{Membaca berita\label{reading_paper}} Dua artikel di \emph{NY Times}
memberitakan studi berikut. Uraian di bawah adalah ringkasan faktual yang
disusun ulang, bukan kutipan dari artikel:
\begin{parts}
\item Artikel berjudul \emph{Risks: Smokers Found More Prone to Dementia} 
melaporkan hal-hal berikut. \footfullcite{news:smokingDementia}
\begin{adjustwidth}{1em}{1em}
{\footnotesize Analisis mencakup 23.123 anggota program kesehatan yang berusia
50--60 tahun ketika secara sukarela menjalani pemeriksaan dan survei perilaku
pada 1978--1985. Setelah 23 tahun, sekitar 25\% peserta tercatat mengalami
demensia; 1.136 kasus adalah penyakit Alzheimer dan 416 kasus adalah demensia
vaskular. Dalam perbandingan yang telah disesuaikan terhadap faktor lain,
perkiraan risiko relatif terhadap bukan perokok adalah 37\% lebih tinggi untuk
satu bungkus per hari, 44\% lebih tinggi untuk satu hingga dua bungkus, dan
sekitar dua kali lipat untuk lebih dari dua bungkus per hari.}
\end{adjustwidth}
Berdasarkan studi ini, dapatkah kita menyimpulkan bahwa merokok menyebabkan 
demensia di kemudian hari? Jelaskan alasan Anda.
\item Artikel lain berjudul \emph{The School Bully Is Sleepy} melaporkan hal-hal 
berikut. \footfullcite{news:bullySleep}
\begin{adjustwidth}{1em}{1em}
{\footnotesize Studi University of Michigan menggabungkan laporan orang tua
tentang kebiasaan tidur anak dengan penilaian orang tua dan guru mengenai
perilaku. Sekitar sepertiga siswa memperoleh penandaan masalah perilaku
mengganggu atau perundungan. Gejala gangguan tidur dilaporkan sekitar dua kali
lebih sering pada anak dalam kelompok bermasalah tersebut dan pada anak yang
diidentifikasi sebagai perundung.}
\end{adjustwidth}
Seorang teman yang membaca artikel itu menyimpulkan bahwa gangguan tidur 
menyebabkan perundungan pada anak sekolah. Apakah pernyataan tersebut dapat 
dibenarkan? Jika tidak, bagaimana sebaiknya Anda menjelaskan kesimpulan yang 
dapat ditarik dari studi ini?
\end{parts}
}{}
}






%%%%%
\section{Eksperimen}
\label{experimentsSection}

%\sectionintro{
Penelitian yang dalam pelaksanaannya peneliti menetapkan perlakuan kepada kasus disebut \termsub{eksperimen}{eksperimen}. Jika penetapan ini melibatkan pengacakan, misalnya menggunakan pelemparan koin untuk menentukan perlakuan yang diterima seorang pasien, penelitian tersebut disebut \term{eksperimen acak}. Eksperimen acak sangat penting ketika kita ingin menunjukkan hubungan sebab-akibat antara dua variabel.
%}\setstretch{1.0}


\subsection{Prinsip perancangan eksperimen}
\label{experimentalDesignPrinciples}

\noindent{}Eksperimen acak pada umumnya dibangun berdasarkan empat prinsip.

\begin{description}
\item[Pengendalian.]
    Peneliti menetapkan perlakuan kepada kasus dan berupaya sebaik mungkin
    untuk \term{mengendalikan} setiap perbedaan lain antarkelompok.\footnote{Konsep
      ini berbeda dari \emph{kelompok kontrol}, yang dibahas pada prinsip
      kedua dan pada Bagian~\ref{biasInHumanExperiments}.}
    Sebagai contoh, ketika pasien meminum obat berbentuk pil, sebagian pasien
    mungkin menelannya hanya dengan seteguk air, sedangkan yang lain dengan
    segelas penuh air. Untuk mengendalikan pengaruh banyaknya air yang diminum,
    dokter dapat meminta semua pasien meminum pil dengan 12 ons cair AS
    (sekitar 355~mL) air.
\item[Pengacakan.] Peneliti mengalokasikan pasien secara acak ke kelompok kontrol
    dan kelompok perlakuan untuk menangani variabel yang tidak dapat dikendalikan. Sebagai
    contoh, pola makan dapat membuat sebagian pasien lebih rentan terhadap
    penyakit daripada pasien lain. Dalam jangka panjang, pengacakan
    menyeimbangkan perbedaan yang terukur maupun yang tidak terukur secara
    rata-rata dan membantu mencegah bias alokasi yang tidak disengaja. Namun, dalam
    satu eksperimen, pengacakan tidak menjamin bahwa kelompok-kelompok yang
    terbentuk akan persis sama.
\item[Replikasi.] Dalam sebuah eksperimen, peneliti melakukan \term{replikasi}
    dengan menerapkan kondisi perlakuan secara independen kepada banyak
    \term{unit eksperimen}. Semakin banyak unit independen yang diamati,
    semakin cermat peneliti dapat mengestimasi pengaruh variabel penjelas
    terhadap variabel respons. Replikasi perlakuan di dalam satu eksperimen
    berbeda dari pengulangan seluruh penelitian oleh tim ilmuwan lain untuk
    memeriksa temuan sebelumnya.

\begin{figure}
  \centering
  \Figure
    [Gambar memperlihatkan tiga tahap dari atas ke bawah. Tahap pertama berupa persegi panjang berisi kisi 54 pasien bernomor 1 sampai 54. Pasien berisiko rendah ditandai dengan lingkaran merah kosong, sedangkan pasien berisiko tinggi ditandai dengan segitiga biru penuh. Tahap kedua memisahkan pasien menjadi blok berisiko rendah di kiri dan blok berisiko tinggi di kanan. Pada tahap ketiga, separuh anggota setiap blok dialokasikan secara acak ke kotak Kontrol dan separuh lainnya ke kotak Perlakuan; warna dan bentuk penanda tetap dipertahankan.]
    {0.82}{figureShowingBlocking}
  \caption{Pemblokan berdasarkan variabel yang menggambarkan risiko pasien.
      Pasien mula-mula dibagi menjadi blok berisiko rendah dan berisiko tinggi;
      kemudian, melalui pengacakan, setiap blok dibagi sama banyak antara
      kelompok kontrol dan kelompok perlakuan. Strategi ini memberi kedua kelompok tersebut
      jumlah pasien berisiko rendah yang sama dan jumlah pasien berisiko
      tinggi yang sama.}
  \label{figureShowingBlocking}
\end{figure}

\item[Pemblokan.] Peneliti kadang-kadang mengetahui atau menduga bahwa variabel
    selain perlakuan turut memengaruhi respons. Dalam keadaan ini, mereka dapat
    lebih dahulu mengelompokkan individu berdasarkan variabel tersebut ke dalam
    \term{blok}, lalu mengalokasikan secara acak kasus di dalam setiap blok ke
    kelompok kontrol dan kelompok perlakuan. Strategi ini disebut \term{pemblokan}. Misalnya, untuk
    meneliti pengaruh obat terhadap serangan jantung, kita dapat lebih dahulu
    membagi pasien menjadi blok berisiko rendah dan berisiko tinggi. Selanjutnya,
    separuh pasien dari setiap blok dialokasikan secara acak ke kelompok kontrol
    dan separuh lainnya ke kelompok perlakuan, seperti pada
    Gambar~\ref{figureShowingBlocking}. Dengan demikian, kelompok kontrol maupun
    kelompok perlakuan memiliki komposisi risiko yang seimbang.
\end{description}

Ketiga prinsip pertama perancangan eksperimen perlu diterapkan dalam setiap
eksperimen. Buku ini menjelaskan metode yang sesuai untuk menganalisis data dari
eksperimen semacam itu. Pemblokan merupakan teknik yang sedikit lebih lanjut;
metode statistik dalam buku ini dapat diperluas untuk menganalisis data yang
dikumpulkan dengan pemblokan.

\subsection{Mengurangi bias dalam eksperimen pada manusia}
\label{biasInHumanExperiments}

Eksperimen acak merupakan tolok ukur utama untuk pengumpulan data, tetapi tidak
selalu menjamin penilaian bebas bias atas hubungan sebab-akibat. Eksperimen pada
manusia merupakan contoh yang baik tentang bagaimana bias dapat muncul tanpa
disengaja. Di sini kita meninjau kembali penelitian yang menggunakan obat baru
untuk menangani pasien serangan jantung. Secara khusus, peneliti ingin mengetahui
apakah obat tersebut mengurangi kematian pasien.

Para peneliti merancang eksperimen acak karena ingin menarik kesimpulan kausal
tentang pengaruh obat. Sukarelawan penelitian\footnote{Subjek manusia sering
disebut \term{pasien}, \term{sukarelawan}, atau \term{peserta penelitian}.}
dialokasikan secara acak ke dua kelompok. Satu kelompok, yaitu \term{kelompok
perlakuan}, menerima obat. Kelompok lainnya, yang disebut \term{kelompok kontrol},
tidak menerima perlakuan obat.

Bayangkan diri Anda sebagai peserta penelitian. Jika masuk kelompok perlakuan,
Anda menerima obat baru yang Anda harapkan akan membantu. Sebaliknya, peserta
di kelompok lain tidak menerima obat dan hanya dapat berharap partisipasinya
tidak meningkatkan risiko kematian. Perbedaan pengalaman ini mengisyaratkan dua
pengaruh: pengaruh yang menjadi perhatian ialah efektivitas obat, sedangkan yang
kedua ialah pengaruh psikologis yang sulit diukur.

Peneliti biasanya tidak berminat pada pengaruh psikologis tersebut karena dapat
membiaskan hasil. Untuk mengatasinya, peneliti berusaha agar pasien tidak
mengetahui kelompoknya. Jika pasien tidak diberi tahu tentang perlakuan yang
diterimanya, eksperimen disebut \term{tersamar}. Namun, ada masalah: pasien yang
sama sekali tidak menerima perlakuan akan tahu bahwa ia berada di kelompok
kontrol. Salah satu solusinya ialah memberi pasien kelompok kontrol perlakuan
semu yang disebut \term{plasebo}. Plasebo yang meyakinkan membantu mempertahankan
penyamaran. Contoh klasiknya ialah pil gula yang dibuat menyerupai pil perlakuan
yang sebenarnya. Kadang-kadang teramati sedikit perbaikan pada pasien yang
menerima plasebo. Fenomena ini sering disebut \term{efek plasebo}; akan tetapi,
pengamatan tersebut tidak dengan sendirinya membuktikan bahwa plasebo menyebabkan
perbaikan, sebab perjalanan alami penyakit dan faktor lain juga dapat berperan.

Bukan hanya pasien yang perlu disamarkan: dokter dan peneliti pun dapat tanpa
sengaja membiaskan eksperimen. Jika mengetahui bahwa seorang pasien menerima
perlakuan sebenarnya, dokter mungkin tanpa sadar memberi pasien itu lebih banyak
perhatian atau perawatan daripada pasien yang menerima plasebo. Untuk mencegah
bias ini, yang dalam beberapa keadaan terbukti memiliki pengaruh terukur,
kebanyakan penelitian modern menggunakan rancangan \term{tersamar ganda}.
Dalam rancangan ini, baik pasien maupun dokter atau peneliti yang berinteraksi
dengan mereka tidak mengetahui siapa yang menerima perlakuan.\footnote{Sebagian
peneliti yang terlibat tetap mengetahui alokasi perlakuan. Namun, mereka tidak
berinteraksi dengan pasien dan tidak memberi tahu tenaga kesehatan yang
disamarkan tentang perlakuan yang diterima setiap pasien.}

\begin{exercisewrap}
\begin{nexercise}
Tinjau kembali penelitian pada Bagian~\ref{basicExampleOfStentsAndStrokes},
tempat peneliti menguji apakah stent efektif mengurangi strok pada pasien
berisiko. Apakah penelitian tersebut merupakan eksperimen? Apakah eksperimen
itu tersamar? Apakah eksperimen itu tersamar ganda?\footnotemark{}
\end{nexercise}
\end{exercisewrap}
\footnotetext{Peneliti mengalokasikan pasien ke kelompok-kelompok eksperimen, jadi penelitian
ini merupakan eksperimen. Namun, pasien dapat membedakan perlakuan yang mereka
terima sehingga eksperimen tidak tersamar. Karena pasien tidak tersamar,
eksperimen tersebut juga tidak mungkin tersamar ganda.}

\begin{exercisewrap}
\begin{nexercise}
\label{gp_sham_surgery}%
Untuk penelitian pada Bagian~\ref{basicExampleOfStentsAndStrokes}, dapatkah
peneliti menggunakan plasebo? Jika dapat, seperti apa bentuk plasebo
tersebut?\footnotemark{}
\end{nexercise}
\end{exercisewrap}
\footnotetext{Secara teknis, peneliti dapat menggunakan \term{bedah semu} agar
pasien tidak mengetahui apakah ia menerima komponen terapeutik. Dalam bedah semu,
pasien menjalani langkah prosedural yang menyerupai operasi, tetapi tidak
menerima perlakuan lengkap. Penyamaran semacam ini memungkinkan peneliti
mengukur respons yang berkaitan dengan plasebo, tetapi tetap memaparkan pasien
pada risiko prosedural tanpa manfaat terapeutik langsung.}

Setelah membaca Latihan Terbimbing~\ref{gp_sham_surgery}, Anda mungkin memiliki
banyak pertanyaan tentang etika penggunaan bedah semu sebagai plasebo.
Pertanyaan serupa berlaku lebih umum ketika perlakuan yang mungkin membantu
tidak diberikan kepada kelompok kontrol. Bedah semu menambahkan risiko dari
prosedur itu sendiri. Menahan perlakuan yang manfaatnya belum terbukti dapat
mempertahankan risiko dasar, sedangkan menahan terapi standar yang sudah terbukti
bermanfaat juga dapat meningkatkan risiko. Karena itu, eksperimen yang etis
memerlukan \emph{ekuipoise klinis}: harus ada ketidakpastian yang sungguh-sungguh
di kalangan ahli mengenai perlakuan mana yang lebih baik.

Eksperimen dan plasebo selalu dapat dilihat dari beberapa sudut pandang, dan
pilihan yang secara etis ``benar'' jarang tampak jelas. Misalnya, etiskah
menggunakan bedah semu jika prosedur tersebut menimbulkan risiko bagi pasien?
Di sisi lain, tanpa pembanding bedah semu, kita mungkin mendorong penggunaan
perlakuan mahal yang sebenarnya tidak efektif. Jika itu terjadi, uang dan
sumber daya lain dapat teralihkan dari perlakuan yang sudah diketahui
bermanfaat. Pada akhirnya, kita menghadapi situasi sulit: perlindungan sempurna
tidak dapat sekaligus diberikan kepada pasien yang menjadi sukarelawan sekarang
dan kepada pasien yang mungkin memperoleh manfaat, atau tidak memperoleh
manfaat, dari perlakuan tersebut di masa mendatang.


{\begingroup\let\clearpageforsection\relax\exercisesheader{}\endgroup

% 29

\eoce{\qt{Kondisi pencahayaan dan kinerja ujian\label{light_exam_performance}} Sebuah studi 
dirancang untuk menguji pengaruh kondisi pencahayaan terhadap kinerja ujian mahasiswa. 
Peneliti menduga bahwa kondisi pencahayaan dapat memberikan pengaruh yang berbeda pada 
laki-laki dan perempuan, sehingga keduanya hendak diwakili sama banyak dalam setiap 
perlakuan. Perlakuannya adalah pencahayaan fluoresen dari atas, pencahayaan kuning dari 
atas, dan tanpa pencahayaan dari atas (hanya menggunakan lampu meja). 
\begin{parts}
\item Apa variabel responsnya?
\item Apa variabel penjelasnya? Apa saja levelnya?
\item Apa variabel pemblokannya? Apa saja levelnya?
\end{parts}
}{}

% 30

\eoce{\qt{Suplemen vitamin\label{vitamin_supplement}}
Untuk menilai efektivitas konsumsi vitamin C dosis tinggi dalam mengurangi durasi selesma, 
para peneliti merekrut 400 sukarelawan sehat dari kalangan staf dan mahasiswa di sebuah 
universitas. Seperempat peserta dialokasikan untuk menerima plasebo, sedangkan peserta 
lainnya dialokasikan secara merata ke tiga perlakuan: vitamin C dosis 1~g, vitamin C dosis 
3~g, atau vitamin C dosis 3~g ditambah zat aditif. Pil tersebut harus mulai diminum ketika 
selesma muncul dan dilanjutkan selama dua hari berikutnya. Semua tablet memiliki tampilan 
dan kemasan yang identik. Perawat yang menyerahkan pil yang diresepkan mengetahui 
perlakuan yang diterima setiap peserta, tetapi 
peneliti yang menilai peserta ketika mereka sakit tidak mengetahuinya. Tidak ditemukan 
perbedaan signifikan di antara keempat kelompok pada ukuran durasi maupun tingkat 
keparahan selesma, dan kelompok plasebo mengalami durasi gejala paling 
singkat.\footfullcite{Audera:2001}
\begin{parts}
\item Apakah studi ini merupakan eksperimen atau studi observasional? Mengapa?
\item Apa variabel penjelas dan variabel respons dalam studi ini?
\item Apakah identitas perlakuan disamarkan dari peserta?
\item Apakah studi ini tersamar ganda?
\item Peserta pada akhirnya dapat memilih apakah akan meminum pil yang diresepkan. Kita 
dapat menduga bahwa tidak semuanya akan mematuhi penugasan dan meminum pilnya. Apakah 
ketidakpatuhan ini dengan sendirinya menciptakan variabel perancu? Jelaskan mengapa alokasi 
awal tetap acak, serta bagaimana ketidakpatuhan dapat mempersulit estimasi efek perlakuan 
yang benar-benar diterima dan membiaskan analisis yang mengelompokkan peserta menurut pil 
yang benar-benar mereka minum.
\end{parts}
}{}

% 31

\eoce{\qt{Kondisi pencahayaan, kebisingan, dan kinerja ujian\label{light_noise_exam_performance}} 
Sebuah studi dirancang untuk menguji pengaruh kondisi pencahayaan dan tingkat kebisingan 
terhadap kinerja ujian mahasiswa. Peneliti menduga bahwa kedua faktor tersebut dapat 
memberikan pengaruh yang berbeda pada laki-laki dan perempuan, sehingga keduanya hendak 
diwakili sama banyak dalam setiap perlakuan. Perlakuan pencahayaan yang dipertimbangkan 
adalah pencahayaan fluoresen dari atas, pencahayaan kuning dari atas, dan tanpa pencahayaan 
dari atas (hanya menggunakan lampu meja). Perlakuan kebisingannya adalah tanpa kebisingan, 
kebisingan konstruksi, dan suara percakapan manusia.
\begin{parts}
\item Apa jenis studi ini?
\item Berapa faktor yang dipertimbangkan dalam studi ini? Sebutkan faktor-faktor tersebut 
dan uraikan levelnya.
\item Apa peran variabel jenis kelamin dalam studi ini?
\end{parts}
}{}

% 32

\eoce{\qt{Musik dan pembelajaran\label{music_learning}} Anda hendak melakukan eksperimen di 
kelas untuk mengetahui apakah mahasiswa belajar lebih baik ketika belajar tanpa musik, 
dengan musik tanpa lirik (instrumental), atau dengan musik berlirik. Uraikan secara 
ringkas sebuah rancangan untuk studi ini.
}{}

% 33

\eoce{\qt{Preferensi minuman bersoda\label{soda_preference}} Anda hendak melakukan eksperimen 
di kelas untuk mengetahui apakah teman-teman sekelas Anda lebih menyukai rasa Coke biasa 
atau Diet Coke. Uraikan secara ringkas sebuah rancangan untuk studi ini.
}{}

% 34

\eoce{\qt{Olahraga dan kesehatan mental\label{exercise_mental_health}} Seorang peneliti tertarik 
pada pengaruh olahraga terhadap kesehatan mental dan mengusulkan studi berikut: Gunakan 
pengambilan sampel acak berstrata untuk memastikan bahwa proporsi kelompok usia 18--30, 
31--40, dan 41--55 tahun dalam sampel mewakili populasi. Selanjutnya, alokasikan secara 
acak separuh subjek dalam setiap kelompok usia untuk berolahraga dua kali seminggu, dan 
instruksikan subjek lainnya agar tidak berolahraga. Lakukan pemeriksaan kesehatan mental 
pada awal dan akhir studi, lalu bandingkan hasilnya.
\begin{parts}
\item Apa jenis studi ini? 
\item Apa kelompok perlakuan dan kelompok kontrol dalam studi ini?
\item Apakah studi ini menerapkan pemblokan? Jika ya, apa variabel pemblokannya?
\item Apakah studi ini menerapkan penyamaran?
\item Jelaskan apakah hasil studi dapat digunakan untuk menetapkan hubungan sebab-akibat 
antara olahraga dan kesehatan mental, serta apakah kesimpulannya dapat digeneralisasikan 
ke populasi secara luas.
\item Andaikan Anda ditugaskan untuk menentukan apakah studi yang diusulkan ini layak 
menerima pendanaan. Apakah Anda memiliki keberatan terhadap usulan studi tersebut?
\end{parts}
}{}
}
